\chapter{\pysv\ 1.6--7}
\setlength{\footmarkwidth}{1.8em}
\pagewiselinenumbers


%[18,25]

\section{Sūtras 1.6 and 7 plus the Bhāṣya up to the definition of \emph{pratyakṣa}}
% This is found on Tm8r3 (right); L8r8 (L8r has 11 lines!)
\mll{BHINT6}
\bht{kāḥ punas tāḥ kliṣṭākliṣṭāḥ pañcaprakārā vṛttaya} ity āha---\edn{don't forget the reference in BSBh 2.4.12: nanv atrāpi śrotrādinirapekṣā bhūtabhaviṣyadādiviṣayā aparā manaso vṛttir astīti samānaḥ pañcasaṃkhyātirekaḥ; evaṃ tarhi 'paramatam apratiṣiddham anumataṃ bhavati' iti nyāyāt ihāpi yogaśāstraprasiddhā manasaḥ pañca vṛttayaḥ parigṛhyante -- 'pramāṇaviparyayavikalpanidrāsmṛtayaḥ' nāma;}
\donote{BHINT6}{kāḥ punas tāḥ kliṣṭākliṣṭāḥ pañcaprakārā vṛttayaḥ}

% ya iti starts Tm8r4
\mll{STR6}%
\str{pramāṇaviparyyayavikalpanidrāsmṛtaya} iti||6||%
\donote{STR6}{\textbf{pramāṇaviparyyaya\- vikalpa\-nidrāsmṛtayaḥ||6||}}

%[18,27]
\vrt{etāvatya eva}{\Tm\Me}{\rdg{etāvaty eva}{L}} vṛttayaḥ| tatra—\edn{should this tatra be part of the sūtra or the bhāṣya?}

\mll{STR7}%
\str{pratyakṣānumānā\vrtms{gamāḥ}{\Tmpc L\Me}{\rdg{°gamā`ḥ'}{\Tm}} pramāṇāni||7||}%
\donote{STR7}{\textbf{pratyakṣānumānāgamāḥ pramāṇāni||7||}}%

\vrt{pramāṇākhyā}{\eme}{\rdg{pramākhyā}{\Tm L\Me}} \mds vṛttis tridhaiva bhidyate| tatra pramā\mde ṇākhyāyāś \vrt{citta\-vṛ\-tteḥ}{\Tm\Lpc\Me}{\rdg{cittavṛtte`ḥ'}{L}} prathamo bhedaḥ %(19,5)
pratyakṣam| atas tallakṣaṇam evābhidhīyate tatpūrvvaka\-tvād bhedadvayasya||
\edn{That all the other pramāṇas presuppose pratyakṣa is not particularly special but it is interesting nonetheless that our author makes an explicit statement.}%

\mll{PRTYDEF}
\bht{indriya\vrtms{pranāḍikay\mds ā}{L\Me}{%
	\rdg{pranāḍikāy\lost{1}}{\Tmac},
	\rdg{pranāḍikay\lost{1}}{\Tmpc}}}---%
indriyam iti \mde \vrt{śrotrādi}{\Tm\Me}{%
	\rdg{śrotrābhi°}{L}}
pañcakam, na karmme\-ndriyam, tasya cittavṛtya\vrtmm{hetuka}{\Tm L}{%
	\rdg{°hetu°}{\Me}}%
\mds tvāt|
\vrtsm{tat}{L}{%
	\rdg{\om}{\Me}}%
kriyārtthaṃ hi tat|
\edn{Perhaps tatkriyārthaṃ tat is a good reading. Cittavṛttyahetuka should refer to śrotrādi (they do not act because of cittavṛtti); "they (śrotrādi) exist for the sake of that (cittavṛtti)" remove "ka"! remove the first tat!}%
tasmād indriya\-śabdasya sā\-māny\mde ārtthatve ’pi \vrt{cittavṛtter}{\Me}{%
	\rdg{ucitavṛtter}{\Tm L}%,
	% \rdg{cittavṛtter}{\Me}%
	}\edn{tu and u are hardly distinguishable. This reading might be justified. However, it would be a little peculiar syntax...}
upa\mds vyācikhyā\mde \vrtmm{sitatvā}{\Tmpc L\Me}{%
	\rdg{°sisatvā°}{\Tmac}}%
c chrotrādi\vrtms{viṣayataiva}{\Tm\Me}{%
	\rdg{°viṣatayaiva}{L}}\edn{The verb upavyākhyā appears in the ChU. Due probably to this, Śaṅkara uses this verb in a technical sense, it seems. Now I have a note on this peculiar verb.}%
| vṛtyartthatvāc ca śrotrādī\mds nām|\edn{again, here appears to be a support that śrotrādi are there for the sake of cittavṛtti}
indriyam eva pranāḍikā dvāram---%
śabdādyā\mde kāra\-vṛttirūpeṇa pariṇamam\mds ānasya \bht{cittasya}|
\vrt{atas}{L}{%
	\rdg{atas tat}{\Me}}
tenendriyadvāre\-ṇa \bht{sāmānya\-vi\-śeṣ\mde ātmakabā\mds hyava\mde stv}ā\mds kāratayā pariṇamamā\mde nam upa\-ra\-jyate|
tasya tad\-\bht{upa\-rāgād} dhetoś cittasya yā % L8v1 starts after this
mudrā\mds pratimudrāvad \bht{vṛttiḥ}| sāmānya\-vi\-śeṣātmaka\-va\mde stūparāge 'pi
\bht{%
	\vrtsm{viśeṣāva\mds dhāraṇa}{L\Me}{%
		\rdg{vi\wmhl{1}eṣāva\lost{4}}{\Tm}}%
	pradhānā%
} saiva \bht{pratyakṣam pramāṇām}||% remove danda before?
\donote{PRTYDEF}{indriyapranāḍikayā cittasya sāmānyaviśeṣātmakabāhyavastūparāgād  viśeṣāvadhāraṇapradhānā vṛttiḥ pratyakṣam pramāṇām||}
%Bhāṣya: indriyapranāḍikayā cittasya sāmānyaviśeṣātmakabāhyavastūparāgād  viśeṣāvadhāraṇapradhānā vṛttiḥ pratyakṣam pramāṇām

\section{\emph{Pratyakṣa} and \emph{viśeṣa}}
\subsection{Why the Bhāṣya says \emph{pratyakṣa} is mainly concerned with difference}
%[19,13]
atha svaviṣayasāmānyaviśeṣāvabhāsakasāma\mde rtthye śrotrādīnāñ cittasya ca
\vrt{kiṃ\-kṛtaṃ}{L\Me}{%
	\rdg{kiṃ\wmhl{1}tam}{\Tm}}
viśeṣāvadhāra\mds ṇaprādhānyam iti---\edn{The choice of the verb avabhās and "ka"?}

ucyate---viśeṣapratibandhatvād\mde\ dhāno\-pādān\mds ādī\mde nāṃ vi\mds
śeṣāvadhāraṇa\vrtms{pra\-dhā\-nety| āha---}{L}{%
	\rdg{pradhānateti (tetyāha)}{\Me}}%
\edn{Something is lost, around here, I think. Something like "(viśeṣāvadhāraṇapradhāneti) yuktam| tadā kiṃ sāmānyan nāvadhāryata ity (āha)" Or simply replace āha with uktam? without the question mark; just divide by a comma? āha part of the previous sentence?}%
na sāmānyaṃ nāvadhāryyate? guṇabhūtan tu gṛhyata eva| yathā nīlaṃ rūpam iti
nīl\mde āvadhāraṇa\vrtms{prādhānye}{\Tmpc L\Me}{%
	\rdg{°pradhānye}{\Tmac}}
saty api rūpam eva \vrt{nīlam}{L\Me}{%
	\rdg{\wmhl{1}lam}{\Tm}}
iti guṇabhūtam a\mds vadhārayati tathā duḥkhaśabdaḥ sukhaśabdo mūḍhaśabda i\mde ti ca| \mds
% Tm8r left is over. This was the last line (line 9).
tathā ca viśeṣāvadhāraṇapradhānā vṛttir iti hy upapadyate||

%[19,20]
kiñ ca saṃśayaviparyyayanimittatvāc ca sāmānyopalabdheḥ| % L8v4 starts after tvā
\edn{"kiñ ca" and "...tvāc ca" are redundant.}%
\vrtsm{sāmānye hy upa}{\Tmac\Tmpc L\Me}{%
	\rdg{sāmānye \ccs hy u\cce `hy u'pa°}{\Tm}}%
labdhe bubhutsitaviśeṣasya saṃśayo jā\mds yate|
% Tm8v1 after jā
samupalabdhasāmānyasmṛter vvā viparyyayaḥ| samadhigataviśeṣasya tu \vrt{saṃśayaviparyayau}{\Me}{\rdg{saṃśayaviparyayā}{L}} na sthitim
% L8v5
upā\-śnuvāte iti viśeṣāva\mde dhāraṇapradh\mds ānatocyate||
\edn{Sāmānya being the origin of saṃśaya and viparyaya could be a unique point of view. At least our author is following something. I should figure this out.}%

%[19,24]
yatrāpi sāmānyam evāvadidhārayiṣitaṃ gaur iti vāśva \vrt{iti vā}{\Me}{%
	\rdg{iti vātra}{L}}%
, tatrāpy anyata\-ra\-parityāgenānya\vrtms{tarāvadhāraṇeti}{\Me}{%
	\rdg{°tarāva\ilg dhāraṇeti}{L}} % L8v6 starts after vi next line
viśeṣāvadhāraṇapradhānataiva| sarvam e\-va hi \vrt{vastu}{\Lpc \Me}{%
	\rdg{vastu \ccs ta\cce}{L}}
vastvantarāpekṣayā sāmānyaṃ viśeṣaś ca||

%[19,27] A new view
\subsection{Criticism of the view that \emph{sāmānya} is unreal (Buddhist)}
yasya punaḥ sāmānyam avastv adhyāropitam evoṣarāmbuvat nendriyaviṣayam iti|
% L8v7 after ta next line
\edn{Certainly a Buddhist but who? Dignāga?}%
\edn{The matter is a bit more complicated. See the Tarkabhāṣā of Mokṣākaragupta toward the end of pratyakṣa and Kajiyama's notes 131ff (section 7.1; pp. 56ff.). According to Kajiyama, Dharmottara and Ratnakīrti contributed to Mokṣākara's saying indriya also takes sāmānya as object but I don't see those two clearly saying it in the passages cited by Kajiyama. The criticism that Buddhist won't acquire sāmānya to form inferences is already in Kumārila.}%
tasya pratiniyata\-sāmānyā\-dhyā\-ropaṇā\-nupapattiḥ| na hy anupala\mde bdhapū\-rvvamu\mds khye \mde sambhava\mds ty adhyāropaṇā, smṛtyasambhavāt|
\edn{The meaning is clear but the use of the compound anupalabdhapūrvamukhya and its association of smṛti might be of interest}%
na kadācid adṛṣṭamukhyodakasya salilā\-dhyā\-ropaṇam ūṣareṣu vidyate||
\edn{Is he talking about some sort of saltlake like places? The example used by the Buddhist means something like positing (salty) water in salt rock desert?}%

%[p]
athāsmarann apy adhyāropayet||
\edn{This, too, must have some background...}%

%[s]
rūpasvalakṣaṇe śabdatvādy apy adhyasyet||
\edn{Again the meaning is more or less clear but apparently this is a very abridged form of a more lengthy discussion.}%

athendriyāntara\mde viṣayatvād \vrt{rūpasvalakṣa\mds ṇe}{\Me}{%
	\rdg{rūpasvalakṣa\lost{1}}{\Tm},
	\rdg{rūpasvalakṣaṇa°}{L}}
\vrt{śabdatvādyadhyāso}{\Lpc\Me}{%
	\rdg{śabdatvā\ccs du\cce dyabhyāso}{L}}
viruddha iti cet---na, avastutve % L8v9 from saty
saty aviśeṣāt| \vrt{viśeṣābhyupagame}{\Lpc\Me}{\rdg{viśeṣā\ccs bhyu\cce bhyupagame}{L}} ca vastutvaprasaṃgaḥ| śa\-bdasvalakṣaṇād iva rūpasvalakṣa\mde ṇasya||
\edn{OK, here is more detail. Not necessarily an abliged discussion. Do we see something like this elsewhere?}%

%[20,6]
ki\mds ñ cānyat---\vrtsm{deśakālānyatva}{\eme}{\rdg{deśakālānya°}{L\Me}}viśi\mde ṣṭa\vrtms{pratyakṣābhyupagame}{\Tm L}{\rdg{pratyakṣā[na]abhyupagame}{\Me}}
\edn{There are only deśakālānyathā or deśakālānyatva in use it should be deśakālādi. See the next sentence. deśakālādyaviśiṣṭa...!}%
samastaloka\mds vya\-va\-hāra\vrtms{vilopa}{\Me(°paḥ)}{%
	\rdg{°vilopā}{L}}
syāt|
%1 ) L8v10 after hī below
na hīdam aham amuṣminn avakāśe kāle cāsminn adrākṣam etasmān nimittād evaṃviśeṣaṇam iti ca \vrt{smṛtyupapattiḥ|
% 2)
na}{\Me}{%
	\rdg{smṛtyupapattin na}{L}}
ca \mde pratyakṣā\vrtms{dṛṣṭe\-ṣu}{L\Me}{\rdg{°dṛ\wmhl{1}eṣu}{\Tm}} viśeṣyavad vi\vrtms{śeṣaṇeṣu}{L\Me}{\rdg{°śeṣaṇe\ccs nu\cce ṣu}{\Tm}} smṛtiḥ saṃbhavati|
% 3)
na ca kalpa\mds nāpoḍhena vya\-va\-hāra\-\vrtms{gocarātītena}{\Me}{%
	\rdg{gocarātī\wmhl{1}e\wmhl{1}a}{L}} % L8v11 at tī
vyavaharamāṇo dṛśyate|
\edn{conclusion p. 273 of Prasannapadā vol. 1.}%
tasmāl lokaprasiddhapratyakṣānu\-saraṇa\-m eva nyāyyam||

\section{\emph{Yogipratyakṣa} and \emph{mānasapratyakṣa}}
% [20,11]
\subsection{Question: Are they separate \emph{pramāṇa}(s) or kinds of \emph{pratyakṣa}?}
nanu ca yoginān ni\mde rvvikalpasamādhijan
\edn{This terminology reflects that of the Nyāyabhāṣya ./6\_sastra/3\_phil/nyaya/nyayabhasya.txt:ubhayathā darśanam, ``asty ātmā'' ity āptopadeśāt pratīyate, tatrānumānam --- ``icchādveṣaprayatnasukhaduḥkhajñānāny ātmano liṅgam'' iti, pratyakṣaṃ --- yuñjānasya yogasamādhijam ``ātmamanasoḥ saṃyogaviśeṣād ātmā pratyakṣa'' iti/ agnir āptopadeśāt pratīyate ``atrāgniḥ''iti, pratyāsīdatā dhūmadarśanenānumīyate, pratyāsannena ca pratyakṣata upalabhyate/ vyavasthā punaḥ --- ``agnihotra juhuyāt svargakāmaḥ'' iti, laukikasya svarge na liṅgadarśanam, na pratyakṣam/ stanayitnuśabde śrūyamāṇe śabdahetor anumānam, tatra na pratyakṣam, nāgamaḥ/ pāṇau pratyakṣata upalabhyamāne nānumānam, nāgama iti/}%
darśanan tasyānevaṃlakṣaṇatvāt tathā sukharāgādi%
\vrtmm{vijñāna\mds sya cānindriya}{L\Me}{%
	\rdg{vijñāna\lost{7}}{\Tm}}%
\vrtmm{pranāḍikā}{L}{%
	\rdg{pranāḍī°}{\Me}}%
\mde pūrvvatvāt \vrt{pramāṇāntaratvam}{\Tm\Me}{\rdg{pramāṇāntaram}{L}} a\-bhy\-upetavyam pratyakṣatvaṃ vā vaktavyam|

\subsection{Answer: \emph{Pratyakṣa}} % Change this at least!
% Consider not emending!
ucyate---\vrt{na| pratyakṣasya}{\conj}{\rdg{na pramāṇasya}{\Tm L}, \rdg{pratyakṣasya (na pramāṇasya)}{\Me}} puruṣapratyayāpekṣatvena \vrt{pari\-hṛta\-tvāt|
na}{\Tmpc L\Me}{\rdg{parihṛtatvena}{\Tmac}} hy apratyayā \vrtsm{vṛttiḥ pra}{\Tm L\Me}{\rdg{vṛttiḥ\ittspc{3}pra°}{\Tm}}tyakṣam|\edn{This part should mean "na pramāṇāntaram". Does this mean what may be called yogipratyakṣa or sukhaduḥkha, etc., are not pramāṇa because they do not [generate] a notion in people? And by puruṣapratyaya, does he mean something like hānopādāna? A tentative translation of this sentence would be "No. For, since pramāṇa is expected to [produce] notions [of hāna and upādān] in man, [they---yogins' darśana and sensations such as sukha-duḥkha] do not qualify [as pramāṇa]." Pratyaya in Śaṅkara is always in the form of a sentence? Not really. Even in the Vivaraṇa itself, there are usages like ahampratyaya or ihapratyaya.}\edn{20161014: The above was probably wrong. The answer "na" is probably OK. This should mean that there is no need to tell that those two are part of pratyakṣa or there is no need to postulate (an)other pramāṇa(s) to incorporate those. "parihṛtatvāt" means that those two possiblities are precluded because of the definition(?) that follows}\edn{Also, this part appears to tell a rathe unique view on pratyakṣa. Do the followings help? BSBh 1.3.33 pratyayāpratyayau hi sadbhāvāsadbhāvayoḥ kāraṇam; nānyārthatvam ananyārthatvaṃ vā; BSBh 2.2.25, 29?}
\edn{Does he mean anumāna, śabda/āptavacana, etc., to be apratyaya?}%
caitanyodaya\vrtms{hetor eva}{\Tm\Lpc\Me}{%
	\rdg{he\ccs tu\cce tor eva}{L}}
sapratya\-yāyāḥ pra\-tya\mds kṣatva\mde m| % L9r2 starts at yāḥ pra
\edn{This sentence basically says the same thing as the previous one.}%
tathā cāha—%
\mll{PRMPHL}%
\bht{phalan tadaviśiṣṭaḥ pauruṣeyaś citta\-vṛtti\-bodha} iti||%
\donote{PRMPHL}{phalan tadaviśiṣṭaḥ pauruṣeyaś citta\-vṛttibodhaḥ}

evañ ca sati sukharāgādijñānasya
\edn{We could emend this to sukharāgādivijñānasya but may not be necessary. The two expressions mean the same thing.}%
\vrt{kliṣṭākliṣṭarūpasya}{\Tm\Me}{\rdg{kliṣṭaṃ kliṣṭaṃ rūpasya}{L}}
\edn{kliṣṭākliṣṭa is an expression to qualify all the vṛttis. See YBh 1.3(?).}%
\vrtsm{tadaviśiṣṭa}{\Tmpc\Me L}{\rdg{tadaviśiṣṭa\ccs viśeṣa\cce}{\Tm}}puruṣapra\-tyayaphalāvasānatvāt pratyakṣatā \vrtsm{siddhai}{L\Me}{\rdg{siddhe°}{\Tm}}va||
\edn{Is the argument as follows?: The Bhāṣya says the result of pratyakṣa is tadaviśiṣṭaḥ pauruṣeyaś cittavṛttibodhaḥ; any cittavṛtti that produces that sort of result qualify as pratyakṣa; the cognition of sukha, rāga, etc., produce such results; therefore they are pratyakṣa? They are not another kind of pramāṇa because...they can be categorized in pratyakṣa?}%
% L9r3 ends with siddhaive


%[21,1]
indriyapranāḍikāgrahaṇaṃ laukika\vrtms{pratyakṣodbhāvanānuvāda}{\Tm}{%
%	\rdg{pratyakṣotbhāvanānuvāda}{\Tm},
	\rdg{pratyakṣotbhāvānuvāda}{L},
	\rdg{pratyakṣodbhavānuvāda}{\Me}}
eva| anyathā hī\-ndriya\vrtms{nirapekṣasya}{\Tmpc L\Me}{%
	\rdg{nirapekṣa\ccs tva\cce sya}{\Tm}}
\vrtsm{yogīśvara}{L\Me}{%
	\rdg{\wmhl{2}gīśvara°}{\Tm}}%
\vrtmm{yor jñāna}{\Tmpc L\Me}{%
	\rdg{yor jjana°}{\Tmac}}%
syāpratyakṣatā syāt|
\edn{So, the first category in quesion, yogipratyakṣa, indeed is pratyakṣa.}%
% Tm9r starts at ādipratibandha
asti hi kleśādi\vrtms{pratibandhābhāve}{\Tm\Lpc\Me}{%
	\rdg{°pratibandhābhāva}{\Lac}}
sarvvārtthena cittasatvena sarvvaviṣayeṇa % L9r4 after satvena sa
yugapatsūkṣma\-vyavahi\-tā\-di\-grahaṇam| tathā cāha---\vrt{tārakaṃ}{\Me}{\rdg{kārakam}{\Tm L}} sarvaviṣayaṃ \vrtsm{sarvvathāviṣayam a}{\Tmpc\Me L}{\rdg{sarvvathāviṣaya`ma'°}{\Tm}}kramañ \vrt{ceti}{L\Me}{\rdg{ce\wmhl{1}}{\Tm}} \vrt{vivekajañ}{L\Me}{\rdg{vi\wmhl{1}ekajaṃ}{\Tm}} jñānam iti| \vrt{tasmān}{L\Me}{\rdg{\wmhl{1}smān}{\Tm}} nāvadhāryyata indriyapranāḍikayai\-veti||
\edn{The use of iti at the end of this paragraph appears interesting.}%

%[21,6] L9r5 starts after bāhya
bāhyavastūparāgād iti viparyyaya\vrtmm{nivṛtyarttham| viśeṣā}{L\Me}{\rdg{°nivṛtyartthaviśeṣā°}{\Tm}}vadhāraṇapradhāna\-tvena saṃśayanivṛttir iti||
\edn{again iti at the end}%

\section{\emph{Pramāṇaphala}}
\subsection{Two views: \emph{pramāṇa} and its \emph{phala} are identical; there is an independent \emph{phala}}
%[21,8]
tatra \vrt{keṣāñ cit}{L\Me(°āṃcit)}{\rdg{ke\wmhl{2}ñcit}{\Tm}} phalan tad eva pramāṇam|\edn{I wonder if the author of the Bhāṣya in fact had Dignāga in mind. The order of the whole commentary (Vivaraṇa) on pratyakṣa reflects that of the Pramāṇasamuccaya. I wonder if this indeed reflects the thought of the Bhāṣyakāra. This seems unlikely. It was the Vivaraṇakāra's idea that somehow the definition in the Bhāṣya should cover yogipratyakṣa and mānasapratyakṣa. The Bhāṣyakāra apparently had no idea his definition would become problematic.}
\edn{Is he basically following the Pramāṇasamuccaya? PS 1.8}%
apareṣān tu pramāṇād \vrt{anyatra}{\Tm L}{\rdg{anyat(tra)}{\Me}} \vrt{pramāṇavyaktyantare}{\Tm L}{\rdg{pramāṇavyaktyantaram(re)}{\Me}}|
\edn{The above sentence looks really weird. I should be looking at the Tantravārttika, Prakaraṇapañcikā, NVTT, NBh, NMañ for hānopāānopekṣā. For pramāṇavyakti, see the TSP, NVTTP.}%
% L9r6 after gha below
\edn{Our author does not like either view.}%
yathā ghaṭavijñānasya pramāṇāsya
\vrt{hānopādānopekṣā}{\Me}{\rdg{hānopādānāpekṣā}{\Tm}, \rdg{hāno\ccs cā\cce pāadānopekṣā}{L}} buddhayaḥ phalam iti
\edn{This is a quote from the Nyāyabhāṣya! No emendation???}%
\vrt{guṇadoṣatadubhayaśūnyatvaviṣayatvād}{\Me(\eme)}{%
  \rdg{guṇadoṣavadubhayaśūnyatra`tva'viṣaya\-tvād}{\Tm},
  \rdg{guṇadoṣavadabhavaśūnyatvaviṣayatvād}{L},
  \rdg{guṇadoṣatadubhaya(vadabhava)śūnyatvaviṣaya\-tvād}{\Me}%
  }
upādānādibuddhīnām
\vrtsm{pramāṇabhūta}{\Tmpc L\Me}{\rdg{brahmāṇaprata°}{\Tmac}}ghaṭavi\mds jñānād
\vrt{a\mde tyantabhedaḥ}{\Lpc\Me}{%
	\rdg{\lost{1}bhyanta\wmhl{1}e\wmhl{2}}{\Tm},
	\rdg{atyanta\ccs meva\cce bhedaḥ}{L}}%
||
\edn{He is in fact following this Naiyāyika idea of hānopādānādibuddhis being the phala of pramāṇa? Most likely not. At the end, he says "nāpi tadvyatiriktapramāṇāntare."}%

\subsection{There is no independent \emph{pramāṇaphala}}
\vrtsm{ataś ca nānya}{\eme}{\rdg{\wmhl{1}taś cānya°}{\Tm}, \rdg{tataś cānya°}{L\Me}}%
% L9r6 after vi below
\vrtms{viṣayasya}{\Tmpc L\Me}{\rdg{°viṣayīsya}{\Tmac}}\edn{This should refer to another reason that is going to be stated why pramāṇa and pramāṇaphala are not two different things. Therefore, tataś ca (for the same reason) does not seem right. L9r6 after vi below}
\vrt{pramāṇasyānyaprameyasya}{}{%
  \rdg{pramāṇasyānyaprameyasya viṣaye}{\Tm L},
  \rdg{pramāṇasya anya(prameyasya)viṣaya°}{\Me}
}
\edn{Could anyaprameyasya an interpolation from a marginal comment?}%
\vrt{pramāṇavyaktyantare}{\Tmpc}{%
  \rdg{pramāṇavyaktyantara°}{\Tmac},
  \rdg{pramāṇe vyaktyantare}{L},
  \rdg{°pramāṇa(ṇe)vyaktyantaraṃ (re)}{\Me}%
} phalam iti|
\edn{The whole of the above is hopelessly corrupt... but at least vyaktyantare phalam or pramāṇavykatyantare phalam appears consistent. So, that reading might be OK.}%
na hi khadiracchedanasya \vrt{khadiropādānan}{\Lpc\Me}{%
	\rdg{khadiropādānān}{\Tmac},
	\rdg{khadiropādā\ccs na\cce nan}{L}} \vrt{tatparityāgo}{\Tm L\Me}{\rdg{tatparityāgo\ittspc{6}}{\Tm}} vā
\edn{The space left in Tm is curious. There could have been something written and erased although doing so without visible traces in a Malayalam manuscript should not be easy because writing literally leaves the surface scratched. A possible candidate that could come there is something like upekṣaṇa.}%
phalaṃ \vrt{bhavituṃ}{\Me}{%
	\rdg{\ilg bhavitu`ṃ'}{\Tm},
	\rdg{bhavatim}{\Lac}}
\edn{I cannot read what the letter before bha is in Tm; Td ignores it.}%
yuktam|
\vrtsm{viyoga}{\Me(\eme)}{%
	\rdg{api yoga°}{\Tmac},
	\rdg{aviyoga°}{\Tmpc L},
	\rdg{vi(api)yoga°}{\Me}%
	}%
sambandhaphalatvāt tyāgo\mds pādānayoḥ|
% L9r8 after tya
\edn{This sentence should explain why tyāga and upādāna are not the result of cutting of khadira tree but the result of ...yogasambandha. This might be simply viyogasambandhaphalatvāt, meaning that a relationship that is detached from khadiracchedana. The decision of whether taking or abondoning the resulting khadira wood comes not from the cutting but from somewhere else.}%
na ca phalasya phalaṃ sambha\mde vati| \mds dvai\mde dhībhāvaphalavyavahitatvāc ca||
\edn{NS/NBh/NV etc., BAUBh, ĀŚV, Bhāmatī, Upad etc. dvaidhībhāva means "making something into two." Something appears to have gone horribly wrong with these paragraphs. Cutting khadira is a Śābarabhāṣya thing. So, the translation might be something like "[Abondoning or taking of the khadira wood] is detached from the result that is making [the tree] into two."}%

% [21,16]
kiñ ca phala\mds sya \mde cet phalam, kriyānapekṣāprasaṅgaḥ| phalārtthaṃ hi
\vrt{kriyā\-pekṣyate}{\Tmpc L\Me}{%
	\rdg{kriyāpekṣyakṣyate}{\Tmac (\emph{kṣya} appearing twice at the end of a line and at the beginnning of the next line)}}%
| tac ce\mds t pha\mde laṃ vināpi kriyayā phalād api
% L9r9 after kriyayā
\edn{phalād api part should mean something like "simply from the [first] result."}%
\vrtsm{sidhyati, ko 'py a}{\conj}{%
	\rdg{sidhyati ko hy a°}{\Tm},
	\rdg{siddhya ko hy a°}{L}
	\rdg{(ato hy a) tato na hy a°}{\Me}}%
nekāy\mds āsā\-pekṣitaniṣpattikriyām anutiṣṭhet|\edn{No need for conjecture???} \vrtsm{sādhanānve\mde ṣaṇā}{\Me}{\rdg{sādhanānveṣaṇaṣa•ā°}{L}}nupapattiś ca, kri\mds yāyā abhā\-vā\mde t||

%[22,4]
nāpi kriyaiva phalam, \vrtsm{kriyānupa}{L\Me}{\rdg{\wmhl{1}yānupa°}{\Tm}}rama\vrtms{prasaṃgāt}{\Tmpc L\Me}{\rdg{°prasaṃgā}{\Tmac}}|
\vrt{phalāvasānā hi}{\eme}{\rdg{\ccs rtthapā\cce `\ilg\ilg'vasānāpi}{\Tm}, \rdg{phalāvasānāpi}{L}, \rdg{phalāvasanāpi [hi]}{\Me}} \mds kriyā| kri\-yāyāś ca
% L9r10 after phala below
phalatve phalā\mde nabhilāṣitva\vrtms{prasaṅgaḥ}{L\Me}{\rdg{°prasaṃ\wmhl{1}ḥ}{\Tm}}|
\mds tasyā duḥkhatvāt||
\edn{Here it is again. The action being pain.}%

\subsection{\emph{Siddhānta}s}
%[22,6]
\subsubsection{No independent \emph{phala} nor are \emph{pramāṇa} and \emph{phala}  identical}
tasmān na pramāṇasvarūpam eva phalam| \vrt{nāpi}{\Me}{\rdg{nādhi}{L}} tadvya\mde ti\mds rikta\vrtms{pramāṇāntara}{L}{\rdg{pramāṇāntaram}{\Me}} ity āha---\bht{phala\mde n tadaviśiṣṭa} ityādi| \vrt{tadaviśiṣṭas}{L\Me}{\rdg{\wmhl{1}daviśiṣṭas}{\Tm}} tayā \vrt{pramāṇākhyayā}{\Tm\Lpc\Me}{\rdg{pramāṇākhyāyā}{\Lac}} vṛ\mds tyā\-viśi\-ṣṭa\-s tulyas tadākāratvena|
% L9r10 after tadākā
\bht{pau\mde ruṣeyaḥ} puruṣasyāyam \mds bodha iti| puruṣasya hi vṛtyākāratā|
\edn{So far, it appears that our author is saying that in pratyakṣa (or in pramāṇa in general?), bodha and pramāṇa (a vṛtti) share the same ākāra, and the bodha and puruṣa share the same ākāra.}%

\subsubsection{The \emph{phala} is not the resulting \emph{dravya}, either}
tasmān \vrt{na kriyājan}{\eme}{\rdg{nākriyājaṃ}{L\Me}} dravyasvarūpamātraṃ phalaṃ|

\subsubsection{The phala is a special state of the \emph{dravya}}

kin tarhi dravyasyai\mde
vā\vrtms{vasthāntaraviśeṣaḥ kriyājaḥ kriyāparyavasānakāle}{\eme}{%
	\rdg{°vasthāntaravi\wmhl{2}ṣa\ccs kriyā\cce\ins(kāryyā)\ine jaḥ \wmhl{1}ryyavasānakālaḥ}{\Tm},
	\rdg{°vasthāntaraviśeṣakriyājākriyāparyāvasānakālaḥ}{L}, %
	\rdg{°vasthāntaraviśeṣaḥ (ṣa)kriyājaḥ(jā)kriyāparyāvasānakālaḥ}{\Me}%
	}
kri\mds yo\vrtmm{parāmopa}{L}{%
	\rdg{°paramopa°}{\Me}}%
patteḥ phalam|\edn{Is this a parallel? ./6\_sastra/3\_phil/vedanta/atmasiddhi.text:yastu sugatamatāvalambī vijñānābhinnahetujatayā tayorapi tadantarbhāvam abhimanyate; kaṇabhakṣapakṣāśrayaṇena vā tayorātmaviśeṣaguṇatvam;tābhyāṃ sukhaduḥkhādhikaraṇaṃ vyācakṣīta; svataḥsukhītyetadvimarśaṃ vātratyam / rāgadveṣādayastu caitanyasyaivāvasthāviśeṣāstadvadeva pratyakṣībhavantīti na tannidarśanenānumānodayaḥ / sukhaprayuktaviṣayīkāracaitanyaṃ rāgaḥ, tadvirodhaprayuktaviṣayīkāraṃ tadeva dveṣaḥ, bhūtaduḥkhajñānena cetaścalanaṃ śokaḥ, āgāmitajjñānena cetaścalanaṃ bhayamityādi lakṣaṇagranthādevāvagantavyamityalaṃ pravistareṇa / Or in BSBh, look for avasthāntara. It is the cause of bhedābhedabhāva Grammarians? About odana and phala} % L9v1 before tteḥ
phalaniṣpattāv eva \mde hi \vrt{kriyoparāma}{L}{%
	\rdg{ṣi`(kri)'yoparāma}{\Tm},
	\rdg{kriyoparama}{\Me}}
\mds upapadyate|
% bodhaḥ should be typeset in bold

\subsubsection{Either \emph{kriyā} or \emph{dravya} is figuratively the \emph{phala} of the other}
evaṃbhūtasya mukhyaphalatve
\vrtms{pradhānakriyāvyavahitaviprakṛṣṭāyāḥ}{\Lpc}{%
	\rdg{pra\ccs mā\cce dhānakriyāvyavahitaviprakṛṣṭāyāḥ}{L},
	\rdg{pradhānakriyāvyavahita(viprakṛṣṭāyāḥ)°}{\Me}}
kriyāyā dravyan dravyasya vā kriyā phalam ity \vrtsm{a\mde nyonya}{L\Me}{%
	\rdg{\lost{1}nyo\ins nya\ine °}{\Tm}}% L9v2 after sthāpa
pratyupasthā\vrtms{pakatvena}{L\Me}{%
	\rdg{°pakatv\wmhl{1}na}{\Tm}}
pra\-dhāna\-phalābhimukhī\vrtms{karaṇāt}{\Tmpc L\Me}{%
	\rdg{karaṇā\ins t\ine}{\Tm}}
pha\mds latvopacāro na virotsyate||
\edn{So, basically the three things are not the phala: kriyā pramāṇa dravya. The last part I cannot say I fully understand but one can figuratively say either dravya or kriyā is the phala? I still don´t understand in what context our author is talking about the phala. Regular kriyā or specifically pramāṇa? Is there the picture of pramāṇa to kriyā to the change in dravya? Looks like pradhānakriyā is pramāṇa and pradhānakriyāvyavahitaviprakṛṣṭakriyā is the secondary;}%

\paragraph{Criticism: the \emph{phala} is only figurative?}
%[22,14] The following criticism appears reasonable. At any rate, what is not pramāṇaphala is more or less over.
nanu ca bhavato ’py upacaritam eva phalam, vṛtyaviśiṣṭāyāḥ
\edn{buddheḥ? The bhāṣya actually reads "pauruṣeyo bodhaḥ." So, it should be masculine. Or he means vṛttyaviśiṣṭāyā avasthāyāḥ? <- That should be it!}%
\edn{Here is a nice one from the Upad: ./upad.txt:atrāha codakaḥ---avagatiḥ pramāṇānāṃ phalaṃ kūṭasthanityātmajyotiḥsvarūpeti ca vipratiṣiddham| ity uktavantam āha---na vipratiṣiddham| kathaṃ tarhi| kūṭasthanityāpi satī pratyakṣādipratyayānte lakṣyate, tādarthyāt| pratyakṣādipratyayasyānityatve 'nityeva bhavati| tena pramāṇānāṃ phalam ity upacaryate|108|}%
phalatvenābhyupagamāt| puruṣasya ca \vrtsm{śuddhatvāt, tadākāratāyā anṛta}{\Me}{%
	\rdg{śuddhatvād akāratāyānṛta°}{L}}%
tvam| na hi sphaṭikamaṇer % L9v3 after ma
a\mde laktakopadhānā\vrtms{kāratā}{L\Me}{\rdg{°k\wmhl{1}ratā}{\Tm}} satyā||
\edn{upādhāna?}%

\paragraph{Answer: yes}
%[22,17]
bāḍham| tathāpi tv etad eva %Tm9v1 starts here
phalam mukhyam, dṛṣṭatvāt, mukhyaphalāntarādarśanāt, sarvvasyaiva kriyākārakajātasyaitatpha\-lārtthatvāt| etadavasānaṃ hi tat sarvvam|
% L9v4
tathā ca pradarśitam||

\paragraph{Criticism: then, the result being more or less erroneous, no reason to pursue \emph{pramāṇa}}
%[22,20]
nanu ca \vrtsm{mithyā}{\Me}{\rdg{mitthyā°}{\Tm L}}bhūtasyaiva mukhyaphalatvaṃ vṛtyāk\mde ārasya,\edn{This notion that pramāṇa is anyway wrong is present in the intro to the BSBh.}
\edn{vṛtyākāra was puruṣa. This sentence does not make much sense.}%
tathā ca tadabhilāṣānupapattiḥ|
\edn{A tentative translation: For [puruṣa] who takes the shape of the [citta]vṛtti, the main result [then] is nothing but unreal. That way, seeking [the main result of pratyakṣa/pramāṇa] will be impossible. doesn't tad refer to puruṣa?}%

\paragraph{Answer: correct}
bāḍham| \vrtsm{mithyā}{\Me}{\rdg{mitthyā°}{\Tm L}}bhūta % Tm9v2
\mds eva vṛtyanugato bhogaḥ| tathā cāha---satvapuruṣayo\mde r atyantāsaṃkīrṇa\mds yoḥ pratyayāviśeṣo bhoga % L9v5 starts at rṇṇa
\edn{YS 3.35}%
 iti| bhā\-ṣyam api---tayoḥ svarūpopalambhe sati kuto bhoga iti||
\edn{What he is doing in these paragraphs is that pramāṇa is essentially mithyā for mixing up puruṣa with cittavṛttis. Pratyakṣa/pramāṇa does not lead to mokṣa. It rather leads to saṃsāra.}%

\paragraph{\emph{samyagdarśana} (rather than \emph{pramāṇa}/\emph{pratyakṣa}) brings about \emph{nivṛtti}}
ata eva \vrt{ca}{\Me}{\rdg{\om}{\Tm L}} samyagdarśanād \mde ātyantikī \vrt{nivṛttir}{L\Me}{\rdg{ni\wmhl{1}ttir}{\Tm}} upapadyate||
\edn{If we adopt the reading na, samyagdarśana means pramāṇa, not vivekakhyāti. But the following suggest ca is better.}%

\paragraph{\emph{Mokṣa} is brought about only by \emph{samyagdarśana} (not by \emph{pramāṇa} that leads to action)}
%[22,25]
yeṣām \vrt{amithyā}{L\Me}{\rdg{amitthyā}{\Tm}} phalan \vrt{te\mds ṣāṃ}{L\Me}{%
	\rdg{te\lost{1}e\lost{1}}{\Tmac}, % L9v6 before vairāgya
	\rdg{te\lost{3}}{\Tmac}}
vairāgyakāraṇābhāvāt, samyagdarśanā\mde nveṣaṇānupapatte\mds r mmokṣābhāvaḥ| na ca kriyāpariniṣpādyo mo\-kṣaḥ, anityatvaprasaṅgāt||
\edn{This is such a typical Śaṅkara thing to say!}%

\subsection{The reason why \emph{puruṣa}’s nature, being aware of \emph{buddhi}, is not established here/Explanation of the intent of the last statement with regard to \emph{pratyakṣa} in the Bhāṣya}
%[23,4]
atha buddhiboddhṛtvam puruṣasya sa\mde tvaṃ cāsiddham iti cet,
% L9v7 after si above
āha---%
\mll{PRTSMVDBU}%
\bht{pratisaṃvedī buddheḥ %Tm9v4
pu\mds ruṣa ity upariṣṭāt pravedayiṣyāma} iti||\mde%
\donote{PRTSMVDBU}{pratisaṃvedī buddheḥ puruṣa ity upariṣṭāt pravedayiṣyāmaḥ}
\edn{This ends the discussion on pramāṇaphala. He does not believe in pramāṇa or pratyakṣa to lead to mokṣa.}%
\edn{The part being referred to in the Bhāṣya is probably not YBh 2.17 but YS 2.20. The Vivaraṇakāra refers to this very passage of YBh 1.7 when commenting on the sūtra (2.20) as the one being referred to be established in the future. YBh 2.20 also has "sa puruṣo buddheḥ pratisaṃvedī." In that sense, whether the reading was pravedayiṣyāmaḥ or upapādayiṣyāmaḥ may be determined by looking at MSS for the part where they record the Vivaraṇa on 2.20. The edition reads upapādayiṣyāmaḥ in YVi 2.20 (p. 189)... So, Td reads upapādayiṣyāma(ḥ) in YVi 2.20 (p. 161/PDF p. 61). Both readings are attested in different manuscripts. Does this mean that our author had access to more than one manuscript of the Yogabhāṣya?}%

\subsubsection{Criticism of the Naiyāyika position of \emph{pratyakṣa}, including their position on \emph{pramāṇaphala}}
% [23,6]
ye \edn{Naiyāyika definition of pratyakṣa; still inside the discussion on pramāṇaphala... Better consult NBh 1.1.4 or thereabout}cāpīndriyārtthasannika\mds rṣasya pratyakṣatvam abhyupayanti,
% L9v8 after kriya
na teṣām api samyag abhyupetam, sannikarṣasya kriyāphalatvāt|
sannikarṣa\mde ś ca saṃyogaḥ|
\vrt{sa ca}{L\Me}{\rdg{\wmhl{1}ca}{\Tm}} kriyājas tair iṣyate|
yato yad upajā\mds yate \vrt{tat tasya}{\Me}{\rdg{tat tasya \ilg}{L}} pramāṇam| \vrt{ta\mde c ca}{L\Me}{\rdg{ta\ittspc{4}c ca}{L}} tasya phalaṃ yad \vrt{upajāyate}{\Tm\Lpc\Me}{\rdg{upajā\ccs\ilg\cce yate}{L}}||
\edn{sannikarṣa = pratyakṣa sannikarṣa <- kriyā; sannikarṣa = kriyāphala sannikarṣa = saṃyoga saṃyoga = kriyāja (kriyāphala) i.e., sannikarṣa/saṃyoga comes from kriyā or kriyāphala, i.e., pratyakṣa is kriyāphala (the result of an action) A is born from B then B is A's pramāṇa; A that is born is the result (phala) of B B -> A then B is pramāṇa and A is phala This is a general principle that both sides should accept but according to the Naiyāyikas saṃyoga = sannikarṣa = pratyakṣa = pramāṇa is a phala of an action not the source (pramāṇa); here is a contradiction}%

% [23,9] L9v9 after iti
phalam api \mds sannikarṣo jñānam prati pramāṇam iti cet---
\edn{Then the Naiyāyikas say sannikarṣa = saṃyoga = pratyakṣa = pramāṇa is a phala (kriyāja) but toward jñāna, it is pramāṇa (source), generator}%

tac ca na,\edn{[Siddhāntin] The key here is the distinction between utpādaka and jñāpaka. Utpādaka should refer to sannikarṣa's being a kāraka} jñ\mde āpakapra\mds māṇa\mde pra\mds stāve kārakapramā\mde ṇopanyāsasyā\vrtms{yuktatvāt}{L\Me}{\rdg{°yu\wmhl{1}\ccs sy\cce āt}{\Tm}}| sannikarṣo hi jñānasyotpāda\mds ko \mde na tasya jñāpakaḥ|
% tasya should be understood as jñānasya; so, "na jñānasya jñāpakaḥ" L9v10 after kāraṇa
prameyādhigamāya hi pra\-tyakṣādīny upa\mds nyastāni, na jñānotpattikāraṇam| \vrt{tatprakriyā\mde yāṃ hy aprakṛtam}{\conj}{%
	\rdg{\lost{4}yāṃ hy aprakṛtaṃ}{\Tm},
	\rdg{tatprakriyāyānyaprakṛtam}{L},
	\rdg{tatprakriyāyām a(yānya)prakṛtaṃ}{\Me}} %
prakriyeta||
\edn{prakriyā is something like a themed discussion? Cf. prakaraṇa? The verb pra-kṛ means something like "to talk about something"?}%

%[23,13]
athā\mds pi syāt\mde ---jñānasyāpi prameyatvāt tac cāpy utpadyamānam pramīyata iti tadutpattikāraṇasya sannikarṣasya jñāpakatvam eveti||
\edn{The point of this opponent's argument is jñānotpattikāraṇa is jñāpaka, and hence sannikarṣa is pramāṇa; no need to distinguish utpādaka and jñāpaka; makes sense, no?}%

tac ca na, jñāna\mds ṃ jñanenaiva jñāyeta\mde, teṣān tathābhyupagamāt|
% L9v11 after jñā of jñāyate above
\edn{Where do the Naiyāyikas teach this? The point should be that one cannot identify jñāpaka and jñānasyotpādaka because only another(?) jñāna is the jñāpaka of jñāna?}%
%
\edn{./6\_sastra/3\_phil/buddh/durvekamisra\_dharmottarapradipa.text:nanu kiṃ ghaḍādidṛṣṭāntabalāt prakāśāntaraprakāśyatā jñānasyāstām, āhosvit pradīpadṛṣṭāntasāmarthyāt prakāśāntaranirapekṣatayā svaprakāśatā? atrocyate | jñānasya svaprakāśarūpatvābhāve ghaṭādeḥ svarūpaprakāśanānupapatteḥ svaprakāśataiveṣṭavyā | atra ca prayogaḥ -- yad avyaktavyaktikaṃ na tad vyaktam yathā kiñcit kadācit kathañcid avyaktavyaktikam | avyaktavyaktikaś cāyaṃ ghaṭādiḥ jñānaparokṣatva iti vyāpakānupalabdhiprasaṅgaḥ | jñānasya jñānāntareṇa vyaktau hetur ayam asiddha iti cen na, nīlajñānodayakāle 'siddhatvād hetor nīlasya parokṣatvaprasañjanāt | na ca bhavatām api mate sarvaṃ vijñānam ekārthasamavāyinā jñānena jñāyate | bubhutsābhāve tadabhāvāt, yathopekṣaṇīyaviṣayā saṃvit | tat upekṣaṇīyam eva tāvad avyaktavyaktikatvād avyaktaṃ prasajyeta | na ceyam ./1\_veda/4\_upa/agamasastra\_w\_sankara.text:sopāyaṃ paramārthatattvaṃ laukikaṃ śuddhalaukikaṃ lokottaraṃ ca krameṇa yena jñānena jñāyate tajjñānam / ./1\_veda/4\_upa/bausbh.text:nanvanyajñānenānyanna jñāyata iti / ./6\_sastra/3\_phil/vedanta/bsbh.text:vyatireke hi satyekavijñānena sarvaṃ vijñāyata itīyaṃ pratijñā hrīyeta  /}%
na ca sannikarṣo jñānam| \vrtsm{jñānotpādaka}{L\Me}{\rdg{jñāno\wmhl{1}ādaka°}{\Tm}}tvenābhyupagamāt||
\edn{So, jñāna, not being a jñāna, it cannot be known}%

% [23,16]
atha \edn{This one sounds alright, again}jñānasya jñāne sannikarṣo \vrt{nimittatvāt pramāṇam}{\Me(conj.)}{\rdg{nimittapramāṇatvam}{L\Tm}} iti \vrt{cet---%
adṛṣṭe}{\Tmpc(ced adṛṣṭe°)\Me}{%
	\rdg{ce`d a'dṛṣṭe°}{\Tm}, %
	\rdg{cedaṣṭe}{L}}%
śvare\-cchāviṣayanimeṣonmeṣādīnām\edn{Perhaps both jñānasya and jñāne are necessary in this context. No need to emend, probably} % L10r1 after ādīnām
api nimi\mds ttatvāt \mde pramāṇatve kaḥ pradveṣaḥ? na ca jñānasya jñeyatva ātmanaś ca svabhāvacidrūpatvānabhyupagame śakyānavasthā parihartum||
\edn{"It is not possible to get rid of the anavasthā if jñāna is to be known and as long as one does not accept ātman being essentially the consciousness."}%

% [23,19]
athātmasamavetamātratayaiva % L10r2 after jñā next
jñānasya jñeyatā| na jñānāntaram apekṣyate yenāti\vrtms{prasajyeteti}{\Tm\Lpc\Me}{%
	\rdg{°prasajyete\ccs te\cce ti}{L}} %
cet---

\vrt{tasyāpy}{\eme}{\rdg{asyāpy}{\Tm L\Me}} anyat pratyakṣalakṣaṇaṃ vācyam| \vrtsm{anindriyasannikarṣādilakṣaṇa}{L\Me}{\rdg{anindriyasa\wmhl{1}ka\wmhl{5}ṇa}{\Tm}}tvād \vrt{ātmani}{L\Me}{\rdg{ā\wmhl{2}}{\Tm}} jñānasamavāyatvasya| sannikarṣasya ca jñānotpatti% This was the end of Tm9v L10r2 starts after jñāno
phalatvāt tadu\-tpattivyavadhāne sati tasya \vrt{jñānasya}{\eme}{\rdg{jñānena}{\Tm L\Me}} \vrt{nānantaryyam}{\Tm\Me}{\rdg{\om}{L}}||
% Was there the end of Tm9v somewhere in this paragraph?

% [23,23]
\vrtsm{tasmāt tadā}{\Tm\Me}{%
	\rdg{\om}{L}}%
nan\vrtms{taryyeṇāpi}{\Tm\Lpc\Me}{%
	\rdg{°taryyaṇāpi}{\Lac}}
na sannikrṣaphalatvaṃ jñānasya| % Were these omissions supplied somewhere in L?
\edn{This ends the discussion on sannikarṣa}%

\paragraph{Examinations of the view that ideas of taking or removing are the result of \emph{pramāṇa}}
tathā hānādibuddhīnām api \vrt{rūpādijñānaphalatvānupapattir uktā}{L\Me}{\rdg{rūpādijñāa\wmhl{8}ā}{\Tm}}|
\edn{Well, the earlier discussion.}%
%[23,25] L10r4 after ca ya
vyabhicārāc ca| \vrt{yathaikasminn}{\Tm\Lpc\Me}{%
	\rdg{yathaivasminn}{\Lac}}
\vrtsm{aṃganādi}{\Tm\Lpc\Me}{%
	\rdg{aṃgānādi°}{\Lac}}%
piṇḍaviṣaye jñāne samāne bahūnām \vrt{utpanne}{L\Me}{%
	\rdg{ulpa\ccs t\cce nn`e'}{\Tm}}
kasyacid upādānabuddhiḥ, anyasya heyabuddhiḥ, kasyacid ubhayābh\mds āvaḥ\mde | tatra hānopādānabuddhir vvyabhicarati|
\vrt{anyataratraikam eva jñānan tad}{L\Me}{% L10r5 before kam eva
	\rdg{anyatara\wmhl{1}ai\wmhl{6}ānan tad}{\Tm}}
\vrt{evobhayaṃ}{L\Me}{%
	\rdg{evo\wmhl{1}yaṃ}{\Tm}}
vyabhicaraty upekṣakasya| yac ca yad vyabhicarati na tat tasya phalam||
% t tasya phalaṃ starts Tm10r3.
\edn{There might be still textual problems but overall the idea is clear.}%

%[24,1]
kiñ ca, hānopādānābhāvamātratvād abhāva evopekṣā| na cābhāvaḥ \mds pramā\-ṇa\-sya va\mde stunaḥ \mds % L10r6 after pha next line
phalam \mde bhavitum arhati| rāgadveṣābhyāñ copādānaparityāga\-bu\-ddhī|
\vrt{tadabhāvāt}{L\Me}{%
	\rdg{tadabhā\wmhl{1}t}{\Tm}}
\vrt{kāraṇābhāve}{\Tm\Lpc\Me}{%
	\rdg{kāraṇābhā\ccs bā\cce ve}{L}}
kāryyābhāva iti pramāṇavijñāne saty % Tm10r4 after sa
api hā\-no\-pādānabuddhyabhāva evopekṣeti
% L10r6 after evope
vyabhicārān na hānādibu\mds ddhiphalatvam||
\edn{"...the notions of hāna, etc., are not phala" Again, the text looks more or less alright. There might be a significance with regard to our author's position on abhāva.... rāga -> upādāna(buddhi); dveṣa -> parityāga(buddhi)/hāna(buddhi)}%

%[24,5]
kiñ ca, rāgadveṣatadabhāvā\mde nām apy upādānādibuddhiphalatve pramāṇatva\-prasaṅgaḥ|
\edn{Err, "There would be the undesirable consequence that rāga, dveṣa, and upekṣā are the pramāṇa with regard to the notions of upādāna, etc., being the result"? The idea seems OK given the diagram above.}%
pauruṣeyacittavṛttibodhaphalatve tu na % Tm10r5 after ki L10r8 after kiñci
kiñcid viru\mds ddhyata iti||
\edn{"But nothing contradicts if [the result/phala is] pauruṣeyacittavṛttibodha."??? iti is again significant here}%

\subsubsection{The problem of \emph{saṃvedyatva} of \emph{cittavṛtti}}
%[24,7] % XYZ
\paragraph{Is \emph{pramāṇa} a \emph{karaṇa}, one of \emph{kāraka}s?}

\mde nanu ca cittavṛtteḥ saṃvedyatvam asiddham| na hi ka\mds raṇam eva karma bhavati| na hi dātram eva lūyate||
\edn{Cittavṛtti and saṃvedyatvam; a new topic So, here cittavṛtti = pramāṇa = karaṇa; cittavṛtteḥ saṃvedyatvam is a problem because of the statement "phalaṃ tad(pramāṇākhyacittavṛtti)aviśiṣṭaḥ pauruṣeyaś cittavṛttibodhaḥ| pratisaṃvedī puruṣa ity upariṣṭāt pravedayiṣyāmaḥ" We have come back to the interpretation of the Bhāṣya For this opponent, the above implies that cittavṛtti is a karaṇa and it cannot be a karman}%

nai\mde ṣa doṣaḥ| kriyākāraka\mds yor bbhinnadha\mde rmmatvāt| kriyā\vrtms{nirvṛttim}{\Me}{\rdg{nivṛttiṃ}{\Tm}}
\edn{kriyānivṛtti or nirvṛtti? Rather nirvṛtti, no? The edition reads that so it is probably my typo when I made the e-text}%
prati sa\-m\-bhū\-ya % L10r9 after sambhūya
kārakāṇi \vrt{vyāpriyante| % Tm10r6 after priya
na \mds parasparaṃ}{L}{%
	\rdg{vyāpriyante na \lost{5}}{\Tm},
	\rdg{vyāpriyante parasparaṃ}{\Me}}
viṣayaviṣayibhāvena|  kri\mde yā tu \vrt{sattā\-mā\-treṇaiva}{\Me}{%
	\rdg{sattāmā\wmhl{1}eṇ\wmhl{2}va}{\Tm},
	\rdg{sattātreṇaiva}{L}}
\mds phalahetur iti karaṇatvam asyām upacaryyate| na
\vrt{kārakavat| kurvvad dhi kārakam| yadi kriyāpi kārakaṃ syā\mde n na}{\conj}{%
	\rdg{[\mist{5}sarvva\mist{2}raka\mist{8}]\ccs sva\cce ān na}{\Tm},
	\rdg{kārakavat kurvvad dhi kāraka yadi kriyādhikākrakaṃ syān na}{L},
	\rdg{kārakavat| kurvad dhi kārayed iti kriyādhikārakaṃ syān na}{\Me}}
bhūtam bhavyāyopadiśyeta|| % L10r10 after bhavyā
\edn{This paragraph is totally obscure... Well, this is from the Śābarabhāṣya! Shows up in the ŚBh, NMañ, Bhāmatī, ŚV commentaries, etc. In essence, the author is saying that pramāṇa is kriyā and therefore not a kāraka and therefore not a karaṇa?}%

%[24,12]
\vrtsm{kiñ cānavasthā}{\Tmpc L\Me}{%
	\rdg{kiñcanavasthā°}{\Tmac}}%
prasaṃga\mds ś ca| % Tm10r7
kriyā cet karaṇaṃ syāt, tannirvvarttyayā \mde
\vrt{cānya\-yā kriyayā}{L\Me}{%
	\rdg{cā\wmhl{5}ā}{\Tm}} %
bhavitavya\mds m| na hy anyāṃ kriyām antareṇa tasyāḥ
\vrt{karaṇatvaṃ}{\eme}{%
	\rdg{kāraṇatvaṃ}{L\Me}} %
ghaṭate| yāpi cānyā sāpi
\vrt{karaṇaṃ}{\eme}{%
%	\rdg{karaṇatvaṃ}{},
	\rdg{kāraṇaṃ}{L\Me}} %
syād anyasyāḥ kriy\mde āyāḥ| sāpy anyasyā ity anavasthā| % L10r11 after s of sāpy
tataś ca kriyāsantānānupara\mds māt % Tm10r8
phalānabhinirvvṛttiḥ|
\vrt{tatrānyā}{L}{%
	\rdg{tatrāntyā}{\Me}}
cet kriyānta\mde reṇa vinā phalahe\mds tur bhavet, prathamasya tathātve kaḥ pradveṣaḥ|
\edn{This sentence does not look good, The last sentence should contribute to the point that kriyā is not a kāraka; and it appears that it is still in the context of anavasthā the reading antya, in contrast to prathama, might be good. These two paragraphs are really hard to figure out}%

\paragraph{\emph{Siddhānta}: \emph{cittavṛtti} is \emph{saṃvedya} as the revealer (\emph{abhivyañjaka})}

cittavṛttes tv abhivyañjakatvāt samvedyatety adoṣaḥ|\edn{Does this part have anything to do with Nyāyasūtra Nyāyabhāṣya NS NBh 1.1.22?} % L10v1 starts after ado
\edn{"tu" is here. As opposed to what? Probably kriyā}%
%[24,18]
dṛśyate cā\mde bhivyañjakā\-nāṃ
\vrtsm{svarūpa}{L\Me}{%
	\rdg{svarū\wmhl{1}°}{\Tm}}%
saṃvedyatayaiva karaṇatvam pra\mds dīpādīnām|
% Tm10r9 after dīpādīnām
\edn{pradīpa, etc., are abhivyañjaka; they are saṃvedya; by being naturally known, they are karaṇa; so is the cittavṛtti...}%
tathā cittavṛtter api|

\subparagraph{Why is the eyesight not grasped while it is an \emph{abhivyañjaka}?}

cakṣurādiṣv abhi\mde vyañjakatve \mds 'py agrahaṇād vyabhicāra
\edn{The sight, etc., do reveal things but they are not grasped, i.e., asaṃvedya, therefore there is a vyabhicāra: not everything that reveals does so by being saṃvedya.}%
iti \vrtsm{cen na, samāna}{\conj}{%
	\rdg{cetanasamāna°}{L}
	\rdg{cet, samāna (cetanasamāna)}{\Me}}%
jātīyasaṃkarāt||

katham?
\edn{The following should explain what he means by samānajātīyasaṃkarāt; why the sight, etc., are not saṃvedya "due to the confusion with things of the same kind"(?)}%

% \subparagraph{According to the theory that sense faculties are material}
%[24,21] L10v2 after yeṣa
yeṣān tāvad bhautikānīndriyāṇi,\edn{position A (Naiyāyikas)} teṣā\mde ñ cakṣuṣas taijasatvād bāhyenālokena \vrtsm{tulya}{\Tmpc L\Me}{\rdg{`tu'lya°}{\Tm}}jātīyena
% Tm10v1 after jā
vyatibhedād gṛhyamāṇasyaivāvivekaḥ| yathā pradīpānāṃ bahūnāṃ prabhāvyatibhedāt kasya kā prabheti % L10v3 afte prabheti
na vivicyate| tathā śrotrādiṣv api||
\edn{This paragraph appears to be relatively well-preserved}%

%[24,24]
% \subparagraph{According to the theory that sense faculties are not material}
yeṣām apy abhautikāni,\edn{Who?}
\vrt{teṣām}{\Me}{%
	\rdg{keṣām}{L}}
apy a\mde bhivyañjakatvena samānajātīyatā||
\edn{Thus far is an explanation why cakṣur etc., are not saṃvedya even though they are abhivyañjaka}%

vṛtter
\vrtsm{apy a\mds vyañjaka}{\conj}{%
	\rdg{apy abhivyañjaka°}{L},
	\rdg{abhivyañjaka°}{\Me}}%
tvād agrahaṇam
\vrt{iti cet, tan na}{L}{%
	\rdg{iti| tan na}{\Me}}%
|| % surviving part of Tm10v2 is below
\vrtsm{prāptacaita\mde nyopagraha}{L\Me}{%
	\rdg{\lost{7}pagrahaṇa°}{\Tm}}%
tay\mds ā
% L10v4 starts after vi of vibhāgya
\edn{The expression prāptacaitanyopagraharūpa is from YBh 2.20, 4.22. This adjective modifies buddhivṛtti The point the S has to prove is that a) cittavṛtti is abhivyañjaka and that b) it is not a grāhya...? It is not grāhya because it is a kriyā...? Shouldn't he be arguing that cittavṛtti is saṃvedya and it is a abhivyañjaka? The opponet here is saying "cittavṛtti is not a vyañjaka and that's why it is not cognized." Our author should counter on both counts, no?}%
\vrt{vyañjakatvād}{\conj}{%
	\rdg{vibhāgyatvād}{L},
	\rdg{'bhibhāvyatvād}{\Me}}
avyavahitatvāc ca| sarvvagatatve ’py ātmanaḥ sūkṣmatvena svacchataratvena ca
\vrtsm{cittenā}{\Me}{%
	\rdg{citte tā}{L}}%
vyavadhānāc ca grahaṇaṃ citta\mde vṛtteḥ||
\edn{Several things to note here. Puruṣa is equated as ātman, it is considered to be sarvagata and sūkṣma. This sentence should be explaninng how the grahaṇa (saṃvid) of cittavṛtti occurs and it appears as though it is not happening.}%

%%%% Conclusion!
\subsection{Conclusion: \emph{pramāṇaphala} is what the Bhāṣya says it is}
tasmāc cittavṛtti\vrtms{pramāṇenaiva}{\Tmpc L\Me}{\rdg{°pramāṇainaiva}{\Tmac}} tadākāratay\mds ā pauruṣeyo ’pi bodhaḥ % L10v5 after bodhaḥ; Tm10v3 after taya
saṃvedya\-māna eva phalam iti siddham ||
\edn{This must be the conclusion of the discussion on pramāṇaphala!}%

\subsubsection{How is \emph{cittavṛtti} object of awareness (\emph{saṃvedya})? Not as \emph{pratyakṣa}}
%[25,4]
yeṣāṃ punaḥ
\vrt{pratyakṣeṇa saṃvedyā}{\eme}{%
	\rdg{pratyakṣeṇāsaṃvedyā}{\all}}
cittavṛttis
\vrtsm{teṣām phala}{\Me}{%
	\rdg{teṣām aphala}{L}}%
grahaṇapramāṇā\-bhāvaḥ|
\edn{Who are these people? What's the point of bringing this up? This discussion in fact is related to the one that starts at XYZ Now I'm pretty sure that the reading should be pratyakṣeṇa saṃvedyā. The difference between the view being attacked and our author's is whether cittavṛtti is saṃvedya by being abhivyañjaka or by being pratyakṣa! Since it is hard to make sense out of this discussion, should we emend the text to "pratyakṣeṇa saṃvedyā"? That would be the Buddhist and Naiyāyika position. The last compound should mean the abhāva of both phalagrahaṇa and pramāṇa Should the last part be pramāṇavyāpārābhāvaḥ? So the whole compound is phalagrahaṇapramāṇavyāpārābhāvaḥ?}%
% This is Tm10v and L10v
phale ca gṛhīte \vrtsm{pramāṇāstitva}{\eme}{%
	\rdg{jñānāstitva}{L},
	\rdg{jñānāsiddhatva}{\Me}}%
\mde syārtthāpattiḥ|
\edn{Is he talking about paracittavṛtti? Most unlikely. Arthāpatti suddenly appearing here is also suspicious. Maybe the argument is if phala is grasped, the existence of pramāṇa/jñāna is inevitable. But...}%
na ca phalaṃ gṛhya\-te % L10v6
\vrtsm{pramāṇavyā\-pārāna}{\Tmpc L\Me}{%
	\rdg{pramāṇavyāpārād a}{\Tmac}}%
bhyupagamāt|
\edn{So, "ca" here means "but," stating what happens if we accept the opponent's position. The phala is not grasped and therefore the knowledge or pramāṇa does not exist. Because the opponent cannot accept the operation of pramāṇa because doing so would result in anavasthā situation.}%
%%% THE FOLLOWING IS IRRELEVANT This statement is not incompatible with the idea "pratyakṣeṇa saṃvedyā" because with that view, the final function of pramāṇa is not specified. cittavṛtti is still known by a pramāṇa but what happens then? A translation could be something like "Also, the effect is not grasped because [according to the theory] they do not accept the function of pramāṇa (for its effect being end of up again as an object of a pramāṇa, i.e., pratyakṣa)."
%%% THE FOLLOWING IS IRRELEVANT The big problem I have is pramāṇavyāpārānabhyupagamāt. Does this refer to a doctorinal position that does not accept the utility of pramāṇa or a consequence of the doctorinal position being discussed? It is hard to conceive anyone who does not accept the working of pramāṇa. It should be then that if we accept certain premises held by the opponent, it becomes inevitable that they also do not accept the utility of pramāṇa.
tadabhyupagame cānavasthāprasaṃgaḥ|
yat phalāvagrahaṇāya pramāṇam abhyupagamyate tasyāpy anyat phalaṃ tatpramāṇasyāpy anyad iti na kvacid avatiṣṭheta||
\edn{tad here is pramāṇavyāpāra tasyāpīti tatpramāṇasyāpi so, pramāṇa P1 -> phala Ph1-> vyāpāra V1 is the basic. If phala is cittavṛtti grasped by pratyakṣa then pramāṇa P1 -> (cittavṛttibodha Ph1 = Ph2 <- pratyakṣa P2) -> vyāpāra V1 then pramāṇa P1 -> ((cittavṛttibodha Ph1 = Ph2 <- pratyakṣa P2) = Ph3 <- pratyakṣa P3) -> vyāpāra V1 then}%
%%% THE FOLLOWING IS IRRELEVANT % From this, it does appear that the second alternative above was the case. The opponent did not say they do not accept the utility of pramāṇa.
%%% THE FOLLOWING IS IRRELEVANT % The logic goes something like, if we accept the premise of the opponent (cittavṛtti is a) saṃvedya or b) asaṃvedya by pratyakṣa),
%%% THE FOLLOWING IS IRRELEVANT % I have the feeling this is OK by our author
%%% THE FOLLOWING IS IRRELEVANT % Oh, this is not about pramāṇaphala being the pauruṣeyaś cittavṛttibodhaḥ but about sukha, dveṣa, etc., being pratyakṣa. Our aurhot is attacking the view that does not accept those as pratyakṣa. However, I have difficulty understanding the relationship between that part (sukha, dveṣa, etc., being pratyakṣa) and the pramāṇaphala being cittavṛttibodha part.
%%% THE FOLLOWING IS IRRELEVANT % The above is probably wrong. This is not about sukhaduḥkhadveṣa, etc., but is about pramāṇa, viparyaya, smṛti, nidrā, etc.
\edn{Is Sureśvara dealing with the same type of person? ./6\_sastra/3\_phil/vedanta/naiskarmyasiddhi\_incomplete.txt:na ca sukhaduḥkhādisaṃbandho .avagatyātmanaḥ pratyakṣādipramāṇair gṛhyate yena virodhaḥ pratyakṣādipramāṇair udchāṭyate| katham| śṛṇu| EVen Dharmakīrti explicitly(?) says: rāgadveṣasukhaduḥkhādīnāṃ sarvacittacaittānām ātmasaṃvedanaṃ pratyakṣam avikalpatvāt | Cf. this from the Brahmasiddhi: tathā hi---pramāṇasya phalam arthāntaram anarthāntaraṃ vā sarvaparīkṣakaiḥ pratijñāyate prajñāyate ca| na ca tad asaṃvedyam, tadasaṃvedyatve sarvāsaṃvedyatvaprasaṅgāt, tatkṛtatvāt saṃvedyabhāvasya bhāvānām| na ca saṃvedyam, phalāntarānupalabdher anavasthāprasaṅgāc ca| tasmāt saṃvedyam, ātmaprakāśatvāt; asaṃvedyaṃ ca, viṣayavat karmabhāvābhāvāt| Maṇḍana's position is that pramāṇaphala is both saṃvedya and asaṃvedya.}%

%[25,8]
atha na
\vrt{\dag pramāṇavyāp\mde āraḥ ta\ldots syāt\dag}{L}{%
	\rdg{\lost{5}āraḥ \ccs ta\cce syāt}{\Tm},
	\rdg{pramāṇavyāpāras tasya}{\Me}}% L10v7 after atha na pramāṇa
||
\edn{The original reading in the common ancestor was "pramāṇavyāpāraḥ tasyāttadā." L actually reads: L10v7 starts vyāpāraḥ ta syāttadāpramāṇābhāvāt phalasyāsatvaṃ prasaktaṃ It is probably reasonable to assume something is missing between the ta and syāt. Since vyāpāraḥ ends with a visarga, there was an awareness that it ended an utterance. Then what would follow is tasya...āt. Then the phrase may be followed by another sentence There is probably a significant portion missing!}%

tadā pramāṇābhāvāt phalasyāsattvaṃ prasaktam| yatra hi sadupalambha\-kāni pramāṇāni na vyāpāraṃ gacchanti, tan nāstīty
\vrt{abhyupagacchāmas teṣām}{L}{%
	\rdg{abhyupagacchāmaḥ (teṣām)}{\Me}}%
||
\edn{tad here must be phala, teṣāṃ pramāṇānām. The basic flow of events appears to be pramāṇa -> (pramāṇa)phala -> (pramāṇa)vyāpāra. The removal of teṣām from the text might be preferable but then where did it come from?}%

%[25,11]
athāgṛhyamāṇam api tad eva jñā\mde nam, pra\mds māṇaṃ phalasyāpīti cet---%
% L10v8 after pramā
\edn{This does sound strange. "phalam api"? tad should be phala. "Even when not being grasped, the same [phala] is the knowledge; the phala, too, has a pramāṇa" or rather the second part should be "the pramāṇa is the evidence (pramāṇa) also for the [existence of] the phala [i.e., there is no need for the phala/jñāna to be grasped. We know that it exists because vyāpāra is being produced]" This could be an answer to the anavasthā problem in the sense there is no need for the phala to be grasped Looks like pramāṇa -> phala = jñāna -> vyāpāra}%
na, a\mde bhivyañjakadharmavyatikramāt|
\edn{Be reminded that abhivyañjakadharma belongs to cittavṛtti according to our author}%
\vrt{na hy abhivyañjakam agṛhyamāṇam}{\Tmpc}{
	\rdg{na hy ativyañjakam agṛhyamāṇam}{\Tmac}, % Tm10v6 starts at ṇam
	\rdg{na hy a`bhivyañjakam agṛyamāṇam'}{L},
	\rdg{na hy abhivyañjakaṃ gṛhyamāṇam}{\Me}}
\edn{If bhi and ti looked similar, this is another indication that the text was from a northern script origin.}%
\vrt{abhi\-vyañjayad}{\Lpc\Me}{%
	\rdg{abhivya[ñja\mist{1}]d}{\Tm},
	\rdg{abhivya\ccs ñjaka\cce yad}{L}}
dṛṣṭam|
\edn{Revealer is always grasped. See above! Reminds me of the lamp being a bad example of. But isn´t this complete opposite of what our author argued for? cittavṛtti is saṃvedya/gṛhya for being abhivyañjaka? This is true if we follow the reading of Lpc (and M, A, and E)}%
na hy \vrt{ulkā}{\Me}{\rdg{utkā}{\Tm L}} giriśikhare ’pi nijan dharmmaṃ atikrāmati||
\edn{Perhaps the most plausible understanding of this example is that fire does what it does with or without anyone watching.}%

%[25,14]
athābhivyañjayad eva phalaṃ praty api pramāṇam iti ce\mde t
\edn{"But the same [pramāṇa] revealing [the object] is the pramāṇa also toward the phala." "abhivyañjayat" should modify "pramāṇam" that should be the pramāṇa for phala, too -> the pramāṇa for the phala was the problem}%
siddhaṃ \vrtsm{tarhīpsita}{\Tmpc L\Me}{\rdg{tarhi smita°}{\Tmac}}m asmāka\mds m|
\vrt{rāgādīn\mde āṃ}{\Me}{% L10v9 after dī
	\rdg{ragādīnāṃ}{L}
}
vṛttiviśeṣaṇatvena grahaṇopapatteḥ|
\edn{Are rāga, etc., here refer to the phala? The word grahaṇa/graha is being used for the phala in this context. Be reminded that rāga and dveṣa give rise to upādāna and hāna/tyāga Need to read this a bit more carefully but looks OK in essence.}%
rakto ’ha\mde n dviṣṭo ’ham
\vrt{iti ca}{\Tm\Lpc}{%
	\rdg{iti ce}{\Lac},
	\rdg{iti}{\Me}}%
| yady
\vrtsm{ahampratyaya}{\Me}{%
	\rdg{ahaṃpratyayi}{\Tm L}}%
viśeṣaṇaṃ,
\vrt{tadāpi}{\Me}{%
	\rdg{tathāpi}{\Tm L}}
\vrtsm{vṛttibhāvabhāvi}{\eme}{%
	\rdg{vṛttibhāvābhāvi}{\Tm L\Me}}%
\mds tvenai\-va grahaṇād vṛttiviśeṣaṇ\mde ānā\mds ṃ
\vrt{rāg\mde ādīnāṃ}{\Tm\Me}{%
	\rdg{\om}{L}}\edn{This is weird. Where did the edition get the reading?}
grahaṇam upapadyate| \vrt{nānyathā}{\Tm L}{% L10v10 after nyatha
	\rdg{a(tenā)nyathā}{\Me}}%
| rūpādi\vrtmm{vat pramānā}{\Tm L}{\rdg{van mānā}{\Me}}ntarāpekṣatvaprasaṅgāt| tataś cānavastheti||
\edn{Here again, iti is used very effectively.}%

%[25,19]
\subsubsection{Final word on \emph{pramāṇaphala}}
\tstwll{PNDT}{tasmāt\ldots upekṣaṇīyam}
tasmāt pramāṇa\vrtmm{phalābhijñānābhi}{L}{%
	\rdg{°phalajñānābhi°}{\Tmac},
	\rdg{°phalānabhijñānām abhi°}{\Me}}%
mānamātram eva kevalaṃ paṇḍitam\vrtms{manyā\-nām}{\Tm\Me}{%
%	\rdg{manyānām}{\Tm\Me},
	\rdg{manyamānām}{L}}
ity upe\mds kṣaṇī\mde yam iti||\edn{Tm has a section marker here. So, this was a major end of a discussion.}
\donote{PNDT}{Cf.\ BSSBh intro.: kathaṃ punar avidyāvadviṣayāṇi pratyakṣādīni pramāṇāni śāstrāṇi ceti|
%
ucyate---dehendriyādiṣv ahaṃmamābhimānarahitasya pramātṛtvānupapattau pramāṇapravṛttyanupapatteḥ| nahīndriyāṇy anupādāya pratyakṣādivyavahāraḥ saṃbhavati| na cādhiṣṭhānam antareṇendriyāṇāṃ vyavahāraḥ saṃbhavati| na cānadhyastātmabhāvena dehena kaścid vyāpriyate| na caitasmin sarvasminn asati asaṅgasyātmanaḥ pramātṛtvam upapadyate| na ca pramātṛtvam antareṇa pramāṇapravṛttir asti| tasmād avidyāvadviṣayāṇy eva pratyakṣādīni pramāṇāni śāstrāṇi ca|
%
paśvādibhiś cāviśeṣāt| yathā hi paśvādayaḥ śabdādibhiḥ śrotrādīnāṃ saṃbandhe sati, śabdādivijñāne pratikūle jāte tato nivartante, anukūle ca pravartante; yathā daṇḍodyatakaraṃ puruṣam abhimukham upalabhya māṃ hantum ayam icchatīti palāyitum ārabhante, haritatṛṇapūrṇapāṇim upalabhya taṃ praty abhimukhībhavanti; evaṃ puruṣā api vyutpannacittāḥ krūradṛṣṭīn ākrośataḥ khaḍgodyatakarān balavata upalabhya tato nivartante, tadviparītān prati pravartante, ataḥ samānaḥ paśvādibhiḥ puruṣāṇāṃ pramāṇaprameyavyavahāraḥ| paśvādīnāṃ ca prasiddho 'vivekapuraḥsaraḥ pratyakṣādivyavahāraḥ, tatsāmānyadarśanād vyutpattimatām api puruṣāṇāṃ pratyakṣādivyavahāras tatkālaḥ samāna iti niścīyate| (BSSBh NSP ed pp. 40–43)}
\edn{Again, the combination of tasmāt and iti So, after yeṣām was the guys who do not accept pramāṇaphala according to the 1952 edition but not according to the 1999 edition...}%
\edn{tatkāla BAUBh yathā vijñātadigvibhāgasyāpyakasmāddigviparyayavibhramaḥ. samyagjñānavato'pi cet pūrvavadviparītapratyaya utpadyate, samyagjñāne'pyavisrambhācchāstrārthavijñānādau pravṛttirasamañjasā syāt, sarvaṃ ca pramāṇamapramāṇaṃ saṃpadyeta, pramāṇāpramāṇayorviśeṣānupapatteḥ. etena samyagjñānānantarameva śarīrapātābhāvaḥ kasmādityetatparihṛtam. jñānotpatteḥ prāk ūrdhvaṃ tatkālajanmāntarasaṃcitānāṃ ca karmaṇāmapravṛttaphalānāṃ vināśaḥ siddho bhavati phalaprāptivighnaniṣedhaśrutereva; 'kṣīyante cāsya karmāṇi' 'tasya tāvadeva ciram' 'sarve pāpmānaḥ pradūyante' 'taṃ viditvā na lipyate karmaṇā pāpakena' 'etamu haivaite na tarataḥ' 'nainaṃ kṛtākṛte tapataḥ' 'etaṁ ha vāva na tapati' 'na bibheti kutaścana' ityādiśrutibhyaśca; 'jñānāgniḥ sarvakarmāṇi bhasmasātkurute' ityādismṛtibhyaśca.}%
\edn{An interesting tidbit from the Pramāṇasamuccaya 1.8 of Dignāga: tām upādāya pramāṇatvam upacaryate nirvyāpāram api sat. tad yathā phalaṃ hetvanurūpam utpadyamānaṃ heturūpaṃ gṛhṇātīty kathyate nirvyāpāram api, tadvad atrāpi. This is about pramāṇaphala. Steinkellner's edition pp. 3–4.}%

\end{document}
%Tm10v
[25,21]
anumānam athocyate 'numeyasya
% Tm: anumānam iti sunocyate a`nu'meyasya
tulyajātīyeṣv anuvṛtta iti| anumeyasya dharmaviśiṣṭasya| tad dharmava%Tmv9/10
ttayā tulyajātīyeṣv anuvṛttaḥ sadbhāvamātreṇa, tadviruddham dhamrmakebhyo vyāvṛttas teṣv asann ity arthaḥ| sambandha iti| kaḥ sambandhaḥ? kayor vā? na hi sambandhināv anāpekṣya sambandhaḥ sambandhatām iyānn asau prasiddha%Tm11r1/p41 of pdf1
vac cocyate ya iti||

[25,26] nanu cānumānasya prakṛtatvāl liṅgaliṅginoḥ sambandhaḥ
prasiddha eva nāgamakatvātn na kevalaḥ sambandho% Tmac sambandhe/Tmpc sambandho
gamaka ucyate liṅgasyaiva sambadho vyāvṛttyanuvṛttyoḥ| kaḥ punar asau? liṅgam eva| na hi liṅgād anya%Tm10/11r2
d anuvṛttaṃ vyāvṛttañ ca ta%ñca Tmac
c ca gamakan nanu liṅgam api naiva gamakaṃ 26,5 liṅgisambandhānapekṣaṃ yadi cānapekṣita%<<yadi cānupekṣā>> yadi cānepekṣita° Tm
liṅgisambandhaṃ gamakaṃ yat kiñcid gamakaṃ syāt, tasmāt liṅgaliṅgisambandho vaktavyaḥ||

[26,7] noktatvād anu%Tm10/11r3
meyasya tulyajātīyeṣv anuvṛttaḥ ity ādinaivoktatvān na hy anumeyenāsaṃ%\mds
baddhaṃ liṅ%\mde
ginā tat sadṛśajātīyeṣv anuvartate%anupravarttate tm
bhinnajātīyebhyo vā nivartate| tasmād ya iti ca prasiddham eva saṃbandham anuvadati

[26,10] tadviṣa%Tm10/11r4
yā evam anvaya% yā eva samanvaya° Tm
vyatirekaviśiṣṭasambandiviṣayā vṛttiḥ %\mds
sāmānyāvadhāraṇapradhānā liṅ%\mde
ga% gi Tm
sāmānyāvadhāraṇapradhānānu% °pradhānā anu° Tm
mānan

[26,12] yathā% tathā tm
deśāntaraprāpter gatimac candratārakañ caitravad iti %Tm10/11r5
yady api la%\mds
kṣaṇā%\mde
bhidhāna%\mds
mā%\mde
treṇa tad abhivyaktis tathāpy udāharaṇapratyudāharaṇa%\mds
pradarśanād atitarām ity udāharaṇa%\mde
pratipādanan tatra gatimac ca%\mds
ndratāraka%\mde
m iti candratārakasya gatimattvam anumeyan tatsamāna%Tm10/11r6
jātī%\mds
yeṣv anuvṛttā deśāntarapr%\mde
āptir gatimatsu caitrā% hard to read Tm
diṣ%\mds
u, bhinnajātīyebhyaś cāgatibhyo vindhyādibhyo vyāvṛttety avyabhica%\mde
rantī %\mds ...
%\mde
% Tm tvañ jātīyakam avinābhāvena sambandhaḥ samba %Tm10/11r7 \mds
% \mde dikāryyakāraṇādibhava%\mds　...
nityaṃ gatimatā sambandhāt||

[26,19] bhavatu tāvad aviruddhānāṃ nityasambandhād deśāntaraprāptyādīnāṃ gamyagamakatā, vadhyaghātakādīnāṃ tu sam%\mde
bandhābhāve katham ity ucyate gamyagamakabhāvbhāvena

[26,21] tatrā%Tm10/11r8\mds
pi nityasambandhād adoṣaḥ| na hy atrāpi prāpt%\mde
ilakṣaṇaḥ sambandho%\mds
'bhipreyate| sāmānyatodṛṣṭā(nu)bhāvaprasaṅgāt| na hi candratārakādiṣu deśāntaraprāpteḥ pr%\mde
āptilakṣāṇasambandhaḥ pratyakṣeṇa gṛhyate agṛhyamā%Tm10/11r9 \mds
ṇaś ced gamakaḥ sarvaṃ sarvasya gamakaṃ syāt| na hi kiṃcit kenacid asambaddham, ākāśādīnāṃ sarvatra bhāvāt||

[26,25] tasmād ganyagamakabuddhyutpattihetur eva sambandho anumānavisaye %\mde
tathā ca badhyaghātakayor api anyatarābhāvada%Tm10/11v1 \mds
rśane 'nyatarabhāvo gamyata ity anumānam eva| yasyāpy ucite deśe abhāvadarśanād anyatra bhāvo gamyate| yathā ghātakadarśanena badhyābhāvo 'numīyate| caitrasya jīvato %\mde Tm10/11v1
gṛhābhāvadarśane bahirbhavaś ceti

[27,1] nanu cābhāvārthā%Tm10/11v3 \mds
pattī ete, nānumāne, (na) lakṣaṇāsaṅgateṅ gavādivat||

[27,2] na ya%\mde
thoktalakṣāno%\mds
papatteḥ badhyaghātakayor virodhalakṣaṇaḥ sambandhaḥ| yatra yatra ghātakabhāvas tatra badhyābhāva evety avinābhāvaḥ%\mde
na hi badhyasyābhāvaṃ ghātakabhāvo%ghātaka(kā)bhāvo E
 vyabhicarati badhyasya vā%Tm10/11v3 \mds
bhāvo ghātakā(ka)bhāvam, agnibhāvam iva dhūmabhāvaḥ||

[27,6] %\mde
kathaṃ punar avagamyate %\mds
ghātakabhāvadarśanad eva tadgocare badhyasyābhāvabuddhir iti ? na punaḥ sadupalambhakapramāṇānāṃ gh%\mde
ātaka iva badhye vyāpārā%vyāpāraḥ Tm
bhāvād ity

[27,8] ucyate ghātakabh%Tm10/11v4 \mds
āve vyāpṛ(pyāvṛ)tapramāṇasyaiva badhyābhāva%\mde
buddhiupapattir dhūmabhāve vy%\mds
āpṛ(vṛ)tatapramāṇasya badhyavi%\mde
śeṣābhāvabuddhyutpattir ddhūmaviśeṣavyāpṛtapramā%Tm10/11v5
ṇasyai%mds
vāgniviśeṣabuddhir i%\mde
ty

[27,12] abhāvapramāṇavādinas tu sarvatra% (tu sa)rvv<sma> Tm
pramāṇābhāvasyāviśiṣṭatvād virodhiviśeṣadarśanānuvidhānaṃ virodhyantarābhāvadarśanasyānupapannaṃ, liṅgaviśeṣānuvudhānam iva liṅgiviśeṣadarśanasya| na hi sadupalambhakapramāṇavyāpārābhāvasya kaścid viśeṣaḥ yena viśeṣabuddhis tad viśeṣam anuvidhīyate||

[27,16] atha jajñāsite viśeṣe sadupalambhakapramāṇaābhāvād iti cet tan nāsataḥ smṛtīcchādisahabhāvāpekṣanupapatteḥ| na hy asat sahāyāpekṣaya kāryaṃ nirvartayad dṛṣṭaṃ śaśa-_Viṣāṇādi| dṛṣṭam iheti cet na vipratipattisthānatve dṛṣṭāntābhāvāt| tulyam iti cet na ghaṭādeḥ sato 'pekṣādarśanāt||

[27,20] na cāpi nirvyāpārasyāvastutvād abhāvasya pramaṇabhāvo yoktaḥ| na hi nirastavyāpārakaṃ kiṃcit kāryārambhakaṃ dṛśyate loke| tasmād vyāpārarūpapramāṇajanitā prameyābhāvabuddhiḥ prameyaviṣayatvat ghaṭādibuddhivat| tathā ca vastutvāt pravṛttinivṛttihetutvāt| vipakṣaḥ śaśaviṣāṇādiḥ| tathābhāva[pramā]pramāṇā(ṇa)bhāvajanitā na bhavati, yathā 'bhihitahetu-dṛṣṭāntadarśanāt(-nena)||

[27,25] ghātakadarśanaṃ badhyābhāvabuddher liṅgaṃ tad bhāvitvād agnyādibuddhinimittadhūmādibuddhivat saṃśayanivartakatvāc ca| badhyābhāvabuddhir vā liṅgajanitā niścayarūpatvād agnibuddhivad iti||

[28,1] iha gaṭo nāstīti kila nānumānaṃ liṅgaliṅgisambandhāgrahaṇāt| na hi yathāgnyādibuddheer dhūmadiliṅgānveṣaṇaṃ tathā iha nāsti ghaṭa ity atra liṅgānveṣaṇam asti| prāg eva liṅgānveṣaṇāt buddhyutpatteḥ||

[28,4] naiṣa doṣaḥ ihapratyayaākārasyaiva liṅgatvāt| na hi prathamam eva liṅgabuddhyupajanane sati liṅgānveṣaṇaṃ bhavati| ihetyākāroparañjitā buddhir liṅgam umuparañjakavyatiriktaghaṭādyabhāvabudder ghātakabuddhivad eva| virodhalakṣaṇa eva cātrāpi sambandhaḥ||

[28,8] yatra punar aviruddhāny uparañjakāni tatrāstitvabuddheḥ etenaivārthāpatter anumānatvaṃ vyākhgyātam avinābhāvena sambandhena| tatrāpi jīvato hi devadattasya gṛhābhāve bahir-bhāvaḥ, yadā gṛhe bhāvaḥ na tadā bahirbhāva ity anyatareṇānumānam | sākṣāt dṛśyamānabahir-bhāvābhyantarābhāvatvaṃ dṛṣṭāntaḥ| vṛttyuparañjakānuparañjakatvena hy abhāvagrahaṇasyābhihito hetuḥ||

[28,13] agnyādāv api dhanādiśaktigrahaṇam anumānam eva| kāryanirvartakayor avinā-bhāvasambandhasyātmani dṛṣṭatvād bhujyādau| dṛśyate hi pratyagātmani śakto 'ham idaṃ kāryaṃ kartum, aśakto 'ham ity ātmavad ātmaviśeṣaṇasya śakter grahaṇam| śaktinirvartyam eva dahanādikāryam apīty anumānasiddhiḥ||

[28,17] nanu cābhāvasya pratyakṣatām upagāyanti kecit| yathā kilādarśakena pradīpādinā sati gṛhyamāṇe vastuni na tad avayavini yad upalabhyate tan nāstīty evaṃ sataḥ prakarśakaṃ pramāṇam asad api prakāśayati| tathā codāharanty anaṅkitaṃ vāsaānayety praviśya bhāṇḍāgāram aṅkitānaṅkitavāsas sannidhāv aṅkitavāsograhaṇenaiva
pratyakṣeṇāṅkitābhāvaviśiṣṭaṃ vāsa upalabhyanayati| tasmāt pratyakṣa evābhāvo lakṣaṇāntarānupapatter iti||

[28,23] tan na buddhyākāreṇārthāvadhāraṇānṛtatvaprasaṅgāt| yadākārā hi buddhis tasya pratyakṣaāvadhriyamāṇātmatābhyupagamyate| yathā nīlaṃ vastu pītam ity atrāpi viśeṣaṇaviśeṣyasambandhapratyayo mithyā prāpnoti| tadyatheha ghaṭo nāstītīhapradeśasya ghaṭābhāvenāsati sambandhe sambandhapratyayah, na hi ghaṭābhāvasyehapradeśasya cā[prāpta]prāptilakṣaṇaḥ samavāyābhidhānovā sambandhas tair abhyyatheha devadatta iti| tasmād abhāva(ve)pratyayavā(yā)n nīlaṃ vastv iti viśeṣaṇaviśeṣyasambandhapratyayo api sambandhapratyayatvān mithyā prasajyeta||

[29,4] yathā cāsty apīhadeśaghaṭābhāvasambandhe sabandhapratītis tathaivānīle api vastuny asaty eva nīlatvasambandge nīlatvasambandhapratyayaḥ prāpnoti, tathā sambandhiṣv api| athāsty eva sambandhaḥ dviṣṭhatā viruddhyeta sambandhasya| abhāvasyāvastutvābhyupagamāt||

[29,7] athāpy ekaśataṃ ṣaṣṭhyarthā iti yāvat ṣaṣthyarthāḥ sambandha iṣyeta, tadā dvividha eva sambandhaḥ samavāyākhyaḥ samyogākhyaś ceti na yujyate||

[29,9] athai(tathai)kaśatasyāpi ṣaṣṭhyarthānāṃ ubhayāntaḥpātitvam ucyeta nābhāve 'anyatarasyāpy asambhavāt||

[29,11] kiñcā(ntva)nyat| iha nabhasi kusumaṃ nāstīti sambandhino nabhaso apy apratyakṣatvāt gaganakusumābhāvasyāpratyakṣatvaṃ
prāpnoti| athāpratyakṣatvam eveti ced āpatitaṃ gaganakusumābhāvavad apratyakṣatvam iha bhū-pradeśe ghaṭābhāvasyāpi| viśeṣo vaktavyaḥ||

[29,14] atha vyomno 'pratyakṣatvād iti tan na puṣpasya patato nabhasy astitvadarśanād vyomāpratyakṣatvam ahetuḥ| yathaiva nipatatkusumasadbhāve vyomāpratyakṣatvam ahetus tathā sadbhāvaviruddhagrahaṇe apy ahetuḥ syāt| athākāśāpratyakṣatvād eveti cet, asambandham api yat kiṃcid vyavahitāpratyakṣatvādi so api hetuḥ syāt| atha nabhaḥ-kusumābhāvaḥ pratyakṣa iti cet indriyārthasannikarṣānarthakyaṃ sarvatra prasajyate| tathā ca saty asannikarṣṭopalabdheḥ sarveṣāṃ sarvajñatvaṃ prasaktam||

[29,20] athānumeyaṃ tatreti cet nabhaḥ kusumābhāvavad iha ghaṭābhāvasyāviśeṣād anumānikatva(kīya)m eva siddham| evaṃ ca saty abhāvapramāṇavādino api vyavahitaviprakṛṣṭodakādāv analādyabhāvasyānumeyatvād iha ghaṭābhāvo api tatsāmānyād anumānagrāhya eveti niścetuṃ yuktam||

[29,24] vindhyāś cāprāpter agatir iti ca prayogāntaram| agatir vindhya iti pratijñānam| aprāpter iti hetuḥ caitravad iti kākalocanavad ubhayatra sambadhyate| gatimac candratārakam ity anena sādharmyeṇa, vindhyāś cāaprāpter ity anena vaidharmyeṇa sambandhaḥ evaṃ ca saty anvayavyatirekā(ka)darśanakṛto na doṣaḥ||

[30,7] atha vaitasmin prayoge vyatirekaviparyayadarśanena doṣābhāva ity upanyastam| lakṣaṇasya ca lakṣyaviśeṣeṇohyatvād adoṣaḥ||

[30,9] āptaḥ parānujighṛkṣāprayukto doṣarahito dṛṣṭānumitārthaḥ| tena dṛṣṭo vā anumito vārthaḥ paratra śrotṛviśeṣe svabodhasaṅkrāntaye svādhigamasaṃkramaṇāya śabdena śabda iti vākyaṃ tato viśiṣṭārthapratipatteḥ| tasmāt tadupadiṣṭād vākyāt| [tadarthaviṣayā] vākyārthaviṣayā yā cittavṛttiḥ sānumānavad eva sāmānyāvadhāraṇapradhānā||

[30,14] ittham āptyādidharmagrahaṇapūrvakā śrotur āgamaḥ na vaktuḥ| dṛṣṭānumitārtho hi vaktā tasya laiṅgakam aindriyakaṃ vā taj-jñānam| śrotur eva tu śravaṇapūrvakaṃ niścitaṃ vijñānaṃ laiṅgikaindriyakavyatiriktam| sa āgamaḥ [na] pratyakṣam indriyā[viṣaya]viṣayatvāt yathā-abhihitaliṅgaliṅgisambandhanirapekṣatvāc ca nānumānam| na hi svargāpūrvadevatādiśabdā(bde ')arthānvitaṃ jñānam āsāditaliṅgaliṅgisambandhapurassaram upajāyate| phalaṃ tu sarvatra
pauruṣeyo vṛttibodha eva||

[30,20] yasyāgamasya vaktāgamākhyavṛttyutpattinimitte śabdoccāraṇāt āgamasya prayogopanyāsena pratipāditaṃ phalaṃ tu pūrvavad eva| yadāpy eka eva prayogas ta(thā)dāpi sādhyasādhanābhāvamātrasya vivakṣitatvān na viparītavyatirekadarśanād ity ucyate||

[30,23] aśraddheyārthavaktrabhihitatvenāgama eva viśiṣyate| katham asāv aśraddheyārtho vaktety ata āha na dṛṣṭo 'numito vārthaḥ tena vaktā| śiṣṭalekaprasiddhānyathā pratipādanāt buddhārhatādiḥ| na hi rāgādidoṣānupahatāntaḥkaraṇaḥ prasiddham anyathā nirvarṇayati| sa āgamaḥ plavata ābhāsī-bhavati||

[30,27] yady api pratyakṣānumānaviṣayābhāsadarśanaṃ svaśabdena na kṛtaṃ tathā apy āgamābhāsopadarśanenaiva pūrvatrāpi yathābhihitalakṣaṇaviparītatvena pratyakṣānumānābhāsādi draṣṭavyam arthād uktam||

[31,1] mūlavaktari tv ādyavaktarīśvare tann nimitto ya āgamaḥ sa viplavanimittābhāvānnirviplavaḥ syāt| yathā-artha(ya)vaktṛvat||

[31,3] upamānaṃ (mānaṃ) śabdapūrvakaṃ tadāgamāntaḥpātitvān na
pramṇāntaram| yathā gaur iva gavaya iti vanecarabhihatād eva vākyād avadhṛtasādṛśya punar vane gavayadarśane api
pūrvaśrutavākyopasthāpitasādṛśyasmaraṇamātram eva nādhikam| sādṛśyaṃ tu vākyād eva pūrvapratipannam||

[31,7] nāpi saṃjñāsaṃjñisambandhapratipattita eva mānaṃ tasyā api tadvākyād evāvagatatvāt| na hi gaur ity etāvatānavagatagavayanāmakasya vanagatagavayadarśino api gavayo 'yam iti saṃjñāsaṃjñisambandhāvabodho bhavati||

[31,10] tasmād gaur iva gavaya iti gavayaśabdo api tadvākye 'pekṣitavyaḥ| tathā ca sati pūrvaśrutād eva sarvaṃ pratītam| athāpi syād ayaṃ(dasya)gavaya iti na pūrvabuddhiviṣayam iti tad api tatrabhavatām api paricodanaparihāram ava[pa]śyatām api
pīḍitakaivale[kaule]yakavṛttikoṇo'vasthānam iva lakṣyate||

[31,14] tathā hi deśakālasannikarṣaviprakarṣādibhedenanekopamānabhedaḥ prasajyeta| tathāstv iti cet tat-pūrvadhiyā eātatvam antar-nihitasamastaviśeṣam eva hi vā neyau pūrvam upapāditam| na hi gaur ity ukte deśakālādiviśeṣābhāve na go-śabdād arthapratipattiḥ||

[31,17] kiṃ cānyat na hi saṃjñāsaṃjñisambandhapratipattir upamānaṃ yena hy apamīyate tadupamānaṃ sādṛśyajñānaṃ hi tat| na hi nāmināmasambandhapratipattyā kiṃcid upamīyata uttarakālatvāt| avadhṛtasādṛśyasya hi paścād ā(sya)yaṃ gavaya ity evaṃ so api nipatati||

[31,21] evam tarhi syād aśabdapūrvakaṃ pramāṇāntaraṃ yathā gavayadarśanād etat sadṛśo gaur iti grāmīṇasya buddhiḥ tad api na
prameyābhāvāt| sati hi prameye sadupalmbhakatvābhimatasya pramāṇasya vyāpāraḥ na hi go-gavayayor vyāvṛttam anuvṛttaṃ ca vastutaḥ sādṛśyaṃ nāmaikam asti||

[31,25] nanu ca sādṛśyabuddhir asty eva nāsau nirālambanā bhavitum arhati| satyam evaṃ tatraiva vicāraṇā kim asau vanayūthādipratyayavat ? atha ghaṭādibuddhivad iti| kim cātaḥ yadi vanapratyayavad bhavet tadā yathaiva viśiṣṭasanniveśānugatā mahīruhā evātmavyatirekeṇāvidyamānavanapratyayahetavas tathaiva go-gavayāvayavā eva dvayadarśanasmṛtyapekṣayā svātmavyatirekeṇāsataḥ sādṛśyasya
pratītau kāraṇam iti prāptam avastutvam||

[32,1] atha punar ghaṭādipratyayavad vastutvaṃ pramāṇatvaṃ ca tadā syād iti||

[32,2] tatra kīdṛśaṃ vastv iti vicāryate 32,3  kiṃ go-gavayayor anuvṛttam ekaṃ sādṛśyaṃ gotvagavayatvābhyām anyat ? te eva vānyonyaviśiṣṭe yadi vā tābhyāṃ bhinne go-gavayādhāre dve sādṛśye parasparānupekḍe, sāpekṣe vā ? 32,6  kiṃ vyaktir eva sānyatarajātiviśiṣṭā vā kiṃ jātiḥ sācānyataravyaktiviśiṣṭā vyaktidvayaviśiṣṭā vā sambandho vaiteṣāṃ yathā-vikalpayogam ity ekaikānyonyaviśiṣṭaviśeṣaṇatvādinā bhūyāṃso 'nye api vikalpā udbhāvayitavyāḥ| tatra yady ekaṃ sādṛśyam ubhayagataṃ tadā nopameyatvam ekatvāt yathā gaur iva gaur iti||

[32,10] atha tenānyataravyaktir upamīyate naitad api svātmagatatvāt gaur iva gotvena| atha tadvatyā vyaktyā vyaktyantaram upamīyate naitad api na hi gotvvaṃ khaṇḍamuṇḍopameyatvakāraṇam| atha dharmāntarāpekṣayā sādṛśyaṃ tadbuddheḥ kāraṇam iti naṣṭaṃ tarhi sādṛśyaṃ sa evāstu dharmaḥ tadbuddheḥ kiṃ te sādṛśyahatakena kalpitena||

[32,14] kiṃ ca sa cāpy apekṣyamāṇo dharmaḥ kim ekaḥ sādṛśyavad ubhayatrātha bhinna ity evam ādivikalpāḥ syur ity asambhavaḥ pūrvavad eva| atha bhinnam\ṃ sādṛśyam ucyate tac ca na bhedād eva na hi bhavati gaur ivāśva iti| yathaivaikaikatvād aghaṭamānatā tathā bahubhir apekṣyamāṇaiḥ| etena jātivyaktisambandhādayaḥ pūrvavikalpitāḥ
pratyuktāḥ||

[32,19] atha buddhyutpattyanyathānupapattyā sādṛśyakalpanā saty evaṃ vanavivarapūrvadakṣiṇottarādibuddhiyogānyathānubhavāt [nyathānupapattyā] vastutvaṃ kiṃ na kalpyate| syād eṣā tava buddhir vṛkṣādaya eva tatra tadbuddhijanmahetava iti| yady evaṃ mithyāpratyayo gavayena sadṛśī gaur iti vanabuddhivat| atha ye vyāvṛttāḥ ka ihopamārthaḥ gaur ivāśva iti| athobhaye api te 'vayavā vyāvṛttanuvṛttātmānas tadbudhdihetavaḥ tathāpi vanapratyayavad eva
prāptam||

[32,25] tasmāt parikṣyamāṇaṃ nātivimardanakṣamam iti ta evāvayavā ubhayatropalabhyamānā bhūyāṃsaḥ sādṛśyabuddheḥ kāraṇaṃ vanabuddher iva
pādapāḥ [iti] nopamānaṃ nāma prmāṇāntaram iti siddham||

[32,27] nāpi premeyadvitvāt pramāṇadvitvakalpanam| kutaḥ pramāṇavaśād arthapratīteḥ| na hy arthavaśāt pramāṇakalpanā| yadi cārthavaśāt
pramāṇakalpanā syāt adhunā 33,4       tan na śakyaṃ jānānām asarvajñatvāt caityavandanāt svargo bhavatīty atra caityavandanasya svargahetutvabuddheḥ pratyakṣanumānāviṣayatvān mithyātvaṃ syāt| tad arthaṃ vā pramāṇāntarakalpanam||

[33,7] atha sā buddhir anmānaphalam iti cet pramāṇād anyatra
phalaprasaṅgaḥ evaṃ canirmokṣaḥ(ka)raṇatrayādivākyārthaviṣyabuddhir apīti||

[33,9] nāpi cātra pramāṇalakṣaṇanirūpaṇatātparyaṃ
prasiddhapramāṇādivṛttinirodhābhyupāyābhipradarśanaparatvāt| ata eva na sūtrakāraḥ sūtrayāñcakāra| bhāsyakāreṇa tu kiñcit
prasiddhānuvādamātram avalambitam iti tad anusareṇena kim api jalpitam iti||6-7||
